%=======================================================
% SOLUCIONES IMPARES. Inserte con \input{Unidad_3_Soluciones_Impares_Para_Insertar}
% Requiere solo el preambulo ya incluido en las notas.
%=======================================================
\clearpage
\noindent{\Large\bfseries Soluciones de los ejercicios impares}\par
\bigskip
\par\bigskip\noindent\textbf{Solución Ejercicio 1}\par\nobreak
Una matriz $m\times n$ tiene $m$ filas y $n$ columnas. Aquí contamos $2$ filas y $3$ columnas, por lo que el orden es $\boxed{2\times3}$.

La notación $a_{ij}$ indica fila $i$ y columna $j$. En la primera fila, segunda columna, encontramos $a_{12}=0$; en la segunda fila, tercera columna, encontramos $a_{23}=5$.

\[
\boxed{a_{12}=0,\qquad a_{23}=5}.
\]

\par\bigskip\noindent\textbf{Solución Ejercicio 3}\par\nobreak
Dos matrices del mismo orden son iguales si y sólo si coinciden sus entradas en las mismas posiciones.

En la posición $(1,1)$:
\[
x - 2=5\quad\Longrightarrow\quad x=7.
\]
En la posición $(1,2)$:
\[
3 y=-6\quad\Longrightarrow\quad y=-2.
\]
En la posición $(2,1)$:
\[
z + 1=4\quad\Longrightarrow\quad z=3.
\]
Sustituimos estos valores en la matriz del primer miembro:
\[
\begin{pmatrix}5&-6\\4&7\end{pmatrix}=\begin{pmatrix}5&-6\\4&7\end{pmatrix}.
\]
Como las matrices coinciden entrada por entrada, la solución es $x=7$, $y=-2$, $z=3$.


\par\bigskip\noindent\textbf{Solución Ejercicio 5}\par\nobreak
Una matriz es simétrica exactamente cuando $A^T=A$.

\[
A_1=\begin{pmatrix}1&2\\2&1\end{pmatrix},\qquad A_1^T=\begin{pmatrix}1&2\\2&1\end{pmatrix}=A_1.
\]
La primera es simétrica.

\[
A_2=\begin{pmatrix}1&0\\3&1\end{pmatrix},\qquad A_2^T=\begin{pmatrix}1&3\\0&1\end{pmatrix}\ne A_2.
\]
La segunda no es simétrica: sus posiciones $(1,2)$ y $(2,1)$ son $0$ y $3$.

\[
I_2^T=\begin{pmatrix}1&0\\0&1\end{pmatrix}=I_2.
\]
La identidad sí es simétrica. \textbf{Respuesta:} primera y tercera matrices.


\par\bigskip\noindent\textbf{Solución Ejercicio 7}\par\nobreak
Las dimensiones son $2\times2$ y $2\times2$; como coinciden las dimensiones interiores (2), el producto es $2\times2$.

La entrada del producto se obtiene multiplicando la fila $i$ del primer factor por la columna $j$ del segundo y sumando los productos:
\[
c_{ij}=\sum_{k=1}^{2} a_{ik}b_{kj}.
\]
\textbf{Fila 1 del producto:}
\[
\begin{aligned}c_{11}&=1\cdot 2+2\cdot 1=4\\c_{12}&=1\cdot 0+2\cdot 3=6\end{aligned}
\]
\textbf{Fila 2 del producto:}
\[
\begin{aligned}c_{21}&=0\cdot 2+1\cdot 1=1\\c_{22}&=0\cdot 0+1\cdot 3=3\end{aligned}
\]
Por tanto,
\[
\boxed{AB=\begin{pmatrix}4&6\\1&3\end{pmatrix}}
\]
\textbf{Invertimos el orden de los factores:} ahora cada fila es de $B$ y cada columna es de $A$.

Las dimensiones son $2\times2$ y $2\times2$; como coinciden las dimensiones interiores (2), el producto es $2\times2$.

La entrada del producto se obtiene multiplicando la fila $i$ del primer factor por la columna $j$ del segundo y sumando los productos:
\[
c_{ij}=\sum_{k=1}^{2} b_{ik}a_{kj}.
\]
\textbf{Fila 1 del producto:}
\[
\begin{aligned}c_{11}&=2\cdot 1+0\cdot 0=2\\c_{12}&=2\cdot 2+0\cdot 1=4\end{aligned}
\]
\textbf{Fila 2 del producto:}
\[
\begin{aligned}c_{21}&=1\cdot 1+3\cdot 0=1\\c_{22}&=1\cdot 2+3\cdot 1=5\end{aligned}
\]
Por tanto,
\[
\boxed{BA=\begin{pmatrix}2&4\\1&5\end{pmatrix}}
\]
En particular, $AB\ne BA$; el producto matricial no es, en general, conmutativo.


\par\bigskip\noindent\textbf{Solución Ejercicio 9}\par\nobreak
Las dimensiones son $2\times3$ y $3\times2$; como coinciden las dimensiones interiores (3), el producto es $2\times2$.

La entrada $c_{ij}$ se obtiene multiplicando la fila $i$ de $A$ por la columna $j$ de $B$ y sumando los productos:
\[
c_{ij}=\sum_{k=1}^{3} a_{ik}b_{kj}.
\]
\textbf{Fila 1 del producto:}
\[
\begin{aligned}c_{11}&=2\cdot 1+(-1)\cdot 0+0\cdot 4=2\\c_{12}&=2\cdot 2+(-1)\cdot (-1)+0\cdot 1=5\end{aligned}
\]
\textbf{Fila 2 del producto:}
\[
\begin{aligned}c_{21}&=1\cdot 1+3\cdot 0+2\cdot 4=9\\c_{22}&=1\cdot 2+3\cdot (-1)+2\cdot 1=1\end{aligned}
\]
Por tanto,
\[
\boxed{AB=\begin{pmatrix}2&5\\9&1\end{pmatrix}}
\]

\par\bigskip\noindent\textbf{Solución Ejercicio 11}\par\nobreak
Sea $A=(a_{ij})$ y $B=(b_{ij})$, ambas de orden $m\times n$. Por definición de suma y transposición, la entrada $(i,j)$ de la matriz transpuesta $(A+B)^T$ es
\[
\bigl((A+B)^T\bigr)_{ij}=(A+B)_{ji}=a_{ji}+b_{ji}.
\]
Por otra parte, la entrada $(i,j)$ de $A^T+B^T$ es
\[
(A^T+B^T)_{ij}=(A^T)_{ij}+(B^T)_{ij}=a_{ji}+b_{ji}.
\]
Ambas matrices tienen orden $n\times m$ y todas sus entradas correspondientes coinciden. Por tanto,
\[
\boxed{(A+B)^T=A^T+B^T}.
\]

\par\bigskip\noindent\textbf{Solución Ejercicio 13}\par\nobreak
Comenzamos con la matriz aumentada $[A\mid I]$ y aplicamos operaciones elementales; cada operación actúa sobre toda la fila.

\[
\left[\begin{array}{rr|rr}1&2&1&0\\3&6&0&1\end{array}\right]
\]
Operación $R_{2}\leftarrow R_{2}-3R_{1}$:
\[
\left[\begin{array}{rr|rr}1&2&1&0\\0&0&-3&1\end{array}\right]
\]
La parte izquierda no se convierte en la identidad: $A$ no es invertible.

También se detecta dependencia de filas: $R_2=3R_1$. El determinante es $1(6)-2(3)=0$; ambas observaciones confirman que la inversa no existe.


\par\bigskip\noindent\textbf{Solución Ejercicio 15}\par\nobreak
Por hipótesis, $A$ es invertible, así que existe $A^{-1}$ y satisface
\[
AA^{-1}=A^{-1}A=I.
\]
Para que una matriz $B$ sea la inversa de $A^{-1}$ debemos comprobar que $B$ multiplicada por $A^{-1}$, en ambos órdenes, produce $I$. La igualdad anterior muestra que $B=A$ cumple las dos condiciones. Por unicidad de la inversa,
\[
\boxed{(A^{-1})^{-1}=A}.
\]

\par\bigskip\noindent\textbf{Solución Ejercicio 17}\par\nobreak
Usamos $\det\begin{psmallmatrix}a&b\\c&d\end{psmallmatrix}=ad-bc$:
\[
\begin{vmatrix}4&-2\\3&5\end{vmatrix}=4\cdot 5-(-2)\cdot 3=20-(-6)=26.
\]

\par\bigskip\noindent\textbf{Solución Ejercicio 19}\par\nobreak
Aplicamos operaciones del tipo $R_i\leftarrow R_i+\lambda R_j$, que \textbf{no cambian} el determinante.
\[
\begin{pmatrix}2&1&0\\4&3&1\\0&2&3\end{pmatrix}\xrightarrow{R_2\leftarrow R_2-2R_1}\begin{pmatrix}2&1&0\\0&1&1\\0&2&3\end{pmatrix}\xrightarrow{R_3\leftarrow R_3-2R_2}\begin{pmatrix}2&1&0\\0&1&1\\0&0&1\end{pmatrix}.
\]
La matriz final es triangular superior. Su determinante es el producto de las entradas diagonales. \[
\boxed{\det(A)=2\cdot1\cdot1=2}.
\]

\par\bigskip\noindent\textbf{Solución Ejercicio 21}\par\nobreak
Observamos que la segunda fila es exactamente el doble de la primera: $R_2=2R_1$. Dos filas linealmente dependientes obligan a que el determinante sea cero.

Equivalentemente, aplicar $R_2\leftarrow R_2-2R_1$ produce una fila nula sin cambiar el determinante:
\[
\begin{pmatrix}1&2&3\\2&4&6\\0&1&1\end{pmatrix}\longrightarrow\begin{pmatrix}1&2&3\\0&0&0\\0&1&1\end{pmatrix},\qquad \boxed{\det(A)=0}.
\]

\par\bigskip\noindent\textbf{Solución Ejercicio 23}\par\nobreak
Es triangular; las entradas bajo (o sobre) la diagonal principal son cero. El determinante es el producto de las entradas de esa diagonal:
\[
\det(A)=1\cdot 3\cdot (-2)\cdot 4=-24.
\]

\par\bigskip\noindent\textbf{Solución Ejercicio 25}\par\nobreak
Consideremos $A=\begin{psmallmatrix}2&1&0\\1&2&1\\0&1&2\end{psmallmatrix}$.

Para obtener la adjunta seguimos el orden: \textbf{menores $\to$ cofactores $\to$ transposición $\to$ inversa}.

Comenzamos calculando el determinante de la matriz.
Elegimos la fila 1, que contiene un cero, para reducir el trabajo. El patrón de signos es $C_{ij}=(-1)^{i+j}M_{ij}$, donde $M_{ij}$ es el determinante que queda al borrar la fila $i$ y la columna $j$.

Los signos de los cofactores alternan según la posición:
\[
\begin{pmatrix}
+&-&+\\
-&+&-\\
+&-&+
\end{pmatrix}.
\]
La entrada de la matriz, el menor y el cofactor son tres objetos distintos; al expandir por una fila se multiplica cada entrada por su cofactor.

Para la entrada $a_{11}=2$, borramos su fila y su columna:
\[
M_{11}=\begin{vmatrix}2&1\\1&2\end{vmatrix}.
\]
\[
M_{11}=2\cdot 2-1\cdot 1=3.
\]
\[
C_{11}=(-1)^{1+1}(3)=3.
\]
Para la entrada $a_{12}=1$, borramos su fila y su columna:
\[
M_{12}=\begin{vmatrix}1&1\\0&2\end{vmatrix}.
\]
\[
M_{12}=1\cdot 2-1\cdot 0=2.
\]
\[
C_{12}=(-1)^{1+2}(2)=-2.
\]
Expandimos: \[
\boxed{\det(A)=2\cdot 3 + 1\cdot (-2)=4}
\]
Calculamos ahora los menores y sus correspondientes cofactores. Calculamos $M_{ij}=\det(A_{ij})$ eliminando la fila $i$ y la columna $j$, y después $C_{ij}=(-1)^{i+j}M_{ij}$.

Fila 1 de la matriz de cofactores:
\[
M_{11}=\begin{vmatrix}2&1\\1&2\end{vmatrix}=2\cdot 2-1\cdot 1=3,\quad C_{11}=(-1)^{1+1}M_{11}=3.
\]
\[
M_{12}=\begin{vmatrix}1&1\\0&2\end{vmatrix}=1\cdot 2-1\cdot 0=2,\quad C_{12}=(-1)^{1+2}M_{12}=-2.
\]
\[
M_{13}=\begin{vmatrix}1&2\\0&1\end{vmatrix}=1\cdot 1-2\cdot 0=1,\quad C_{13}=(-1)^{1+3}M_{13}=1.
\]
Fila 2 de la matriz de cofactores:
\[
M_{21}=\begin{vmatrix}1&0\\1&2\end{vmatrix}=1\cdot 2-0\cdot 1=2,\quad C_{21}=(-1)^{2+1}M_{21}=-2.
\]
\[
M_{22}=\begin{vmatrix}2&0\\0&2\end{vmatrix}=2\cdot 2-0\cdot 0=4,\quad C_{22}=(-1)^{2+2}M_{22}=4.
\]
\[
M_{23}=\begin{vmatrix}2&1\\0&1\end{vmatrix}=2\cdot 1-1\cdot 0=2,\quad C_{23}=(-1)^{2+3}M_{23}=-2.
\]
Fila 3 de la matriz de cofactores:
\[
M_{31}=\begin{vmatrix}1&0\\2&1\end{vmatrix}=1\cdot 1-0\cdot 2=1,\quad C_{31}=(-1)^{3+1}M_{31}=1.
\]
\[
M_{32}=\begin{vmatrix}2&0\\1&1\end{vmatrix}=2\cdot 1-0\cdot 1=2,\quad C_{32}=(-1)^{3+2}M_{32}=-2.
\]
\[
M_{33}=\begin{vmatrix}2&1\\1&2\end{vmatrix}=2\cdot 2-1\cdot 1=3,\quad C_{33}=(-1)^{3+3}M_{33}=3.
\]
Reunimos los cofactores en una matriz y la transponemos para obtener la adjunta.
\[
\mathcal{C}(A)=\begin{pmatrix}3&-2&1\\-2&4&-2\\1&-2&3\end{pmatrix},\qquad\adj(A)=\mathcal{C}(A)^T=\begin{pmatrix}3&-2&1\\-2&4&-2\\1&-2&3\end{pmatrix}.
\]
Como el determinante es distinto de cero, obtenemos la inversa mediante la adjunta.
\[
A^{-1}=\frac1{\det(A)}\adj(A)=\frac1{4}\begin{pmatrix}3&-2&1\\-2&4&-2\\1&-2&3\end{pmatrix}=\begin{pmatrix}\frac{3}{4}&- \frac{1}{2}&\frac{1}{4}\\- \frac{1}{2}&1&- \frac{1}{2}\\\frac{1}{4}&- \frac{1}{2}&\frac{3}{4}\end{pmatrix}.
\]
Como verificación, debe cumplirse $A\adj(A)=\det(A)I$.
\[
\begin{pmatrix}2&1&0\\1&2&1\\0&1&2\end{pmatrix}\begin{pmatrix}3&-2&1\\-2&4&-2\\1&-2&3\end{pmatrix}=\begin{pmatrix}4&0&0\\0&4&0\\0&0&4\end{pmatrix}=4I_{3},
\]
\[
AA^{-1}=\begin{pmatrix}1&0&0\\0&1&0\\0&0&1\end{pmatrix}=I_{3}.
\]

\par\bigskip\noindent\textbf{Solución Ejercicio 27}\par\nobreak
Aplicamos operaciones elementales de fila, que no modifican el rango ni el conjunto de soluciones del sistema homogéneo.

\[
\begin{pmatrix}1&2&1\\0&1&1\end{pmatrix}
\]
Operación $R_{1}\leftarrow R_{1}-2R_{2}$:
\[
\begin{pmatrix}1&0&-1\\0&1&1\end{pmatrix}
\]
Hay pivotes en las columnas 1, 2; por tanto, $\boxed{\rang(A)=2}$.

Resolvemos $A\vect x=\vect0$; la matriz reducida representa el mismo sistema homogéneo.

Las variables libres son $x_{3}$. Sean $x_{3}=t_{1}$, con $t_i\in\mathbb R$.

\[
\begin{aligned}x_{1}=t_{1}\\x_{2}=- t_{1}\end{aligned}
\]
\[
\vect x=t_{1}\begin{pmatrix}1\\-1\\1\end{pmatrix}
\]
\[
\boxed{\Ker(A)=\operatorname{span}\left\{\begin{pmatrix}1\\-1\\1\end{pmatrix}\right\}}
\]
Para verificar el resultado, cada vector generador produce el vector cero al multiplicarlo por $A$:
\[
A\begin{pmatrix}1\\-1\\1\end{pmatrix}=\begin{pmatrix}0\\0\end{pmatrix}.
\]
Además, $\rang(A)+\dim\Ker(A)=2+1=3$, igual al número de columnas.


\par\bigskip\noindent\textbf{Solución Ejercicio 29}\par\nobreak
El sistema $A\vect x=\vect b$ equivale a las dos ecuaciones
\[
\begin{cases}x_1+x_2=2,\\x_1+x_2=2.\end{cases}
\]
Las ecuaciones son idénticas: hay una variable libre. Tomamos $x_2=t$ y despejamos $x_1=2-t$.
\[
\boxed{\vect x=\begin{pmatrix}2-t\\t\end{pmatrix}=\begin{pmatrix}2\\0\end{pmatrix}+t\begin{pmatrix}-1\\1\end{pmatrix},\quad t\in\mathbb R}.
\]
Comprobamos:
\[
\begin{pmatrix}1&1\\1&1\end{pmatrix}\begin{pmatrix}2-t\\t\end{pmatrix}=\begin{pmatrix}2\\2\end{pmatrix}.
\]
La matriz tiene dos filas iguales, de modo que $\det(A)=1-1=0$ y no existe $A^{-1}$. No podemos aplicar la fórmula $A^{-1}\vect b$.


\par\bigskip\noindent\textbf{Solución Ejercicio 31}\par\nobreak
\textbf{Método (a): eliminación.} Formamos la matriz aumentada del sistema:
\[
\left[\begin{array}{rr|r}2&1&7\\1&1&4\end{array}\right]\xrightarrow{R_2\leftarrow2R_2-R_1}\left[\begin{array}{rr|r}2&1&7\\0&1&1\end{array}\right].
\]
La segunda ecuación dice $x_2=1$; de la primera, $2x_1+1=7$, y así $x_1=3$.

\textbf{Método (b): matriz inversa.} Calculamos el determinante y la inversa:
\[
\det(A)=2(1)-1(1)=1,\qquad A^{-1}=\begin{pmatrix}1&-1\\-1&2\end{pmatrix}.
\]
Ahora multiplicamos $\vect x=A^{-1}\vect b$:
\[
\vect x=\begin{pmatrix}1&-1\\-1&2\end{pmatrix}\begin{pmatrix}7\\4\end{pmatrix}=\begin{pmatrix}3\\1\end{pmatrix}.
\]
La comprobación en el sistema original es
\[
\begin{pmatrix}2&1\\1&1\end{pmatrix}\begin{pmatrix}3\\1\end{pmatrix}=\begin{pmatrix}7\\4\end{pmatrix},\qquad\boxed{(x_1,x_2)=(3,1)}.
\]

\par\bigskip\noindent\textbf{Solución Ejercicio 33}\par\nobreak
Aplicamos operaciones elementales de fila, que no modifican el rango ni el conjunto de soluciones del sistema homogéneo.

\[
\begin{pmatrix}1&1&0\\0&1&1\\1&2&1\end{pmatrix}
\]
Operación $R_{3}\leftarrow R_{3}-1R_{1}$:
\[
\begin{pmatrix}1&1&0\\0&1&1\\0&1&1\end{pmatrix}
\]
Operación $R_{1}\leftarrow R_{1}-1R_{2}$:
\[
\begin{pmatrix}1&0&-1\\0&1&1\\0&1&1\end{pmatrix}
\]
Operación $R_{3}\leftarrow R_{3}-1R_{2}$:
\[
\begin{pmatrix}1&0&-1\\0&1&1\\0&0&0\end{pmatrix}
\]
Hay pivotes en las columnas 1, 2; por tanto, $\boxed{\rang(A)=2}$.

Resolvemos $A\vect x=\vect0$; la matriz reducida representa el mismo sistema homogéneo.

Las variables libres son $x_{3}$. Sean $x_{3}=t_{1}$, con $t_i\in\mathbb R$.

\[
\begin{aligned}x_{1}=t_{1}\\x_{2}=- t_{1}\end{aligned}
\]
\[
\vect x=t_{1}\begin{pmatrix}1\\-1\\1\end{pmatrix}
\]
\[
\boxed{\Ker(A)=\operatorname{span}\left\{\begin{pmatrix}1\\-1\\1\end{pmatrix}\right\}}
\]
Para verificar el resultado, cada vector generador produce el vector cero al multiplicarlo por $A$:
\[
A\begin{pmatrix}1\\-1\\1\end{pmatrix}=\begin{pmatrix}0\\0\\0\end{pmatrix}.
\]
Además, $\rang(A)+\dim\Ker(A)=2+1=3$, igual al número de columnas.


\par\bigskip\noindent\textbf{Solución Ejercicio 35}\par\nobreak
Contamos 1 filas y 4 columnas, así que el orden es $1\times4$. El número de entradas es $1\cdot 4=4$.

La entrada $a_{12}$ está en fila 1, columna 2. Allí aparece $0$, luego $\boxed{a_{12}=0}$.

La entrada $a_{13}$ está en fila 1, columna 3. Allí aparece $2$, luego $\boxed{a_{13}=2}$.


\par\bigskip\noindent\textbf{Solución Ejercicio 37}\par\nobreak
Contamos 3 filas y 2 columnas, así que el orden es $3\times2$. El número de entradas es $3\cdot 2=6$.

La entrada $a_{12}$ está en fila 1, columna 2. Allí aparece $1$, luego $\boxed{a_{12}=1}$.

La entrada $a_{12}$ está en fila 1, columna 2. Allí aparece $1$, luego $\boxed{a_{12}=1}$.

\emph{Nota editorial:} el enunciado repite la misma entrada; la respondemos en ambas apariciones sin modificar el problema.


\par\bigskip\noindent\textbf{Solución Ejercicio 39}\par\nobreak
Contamos 4 filas y 3 columnas, así que el orden es $4\times3$. El número de entradas es $4\cdot 3=12$.

La entrada $a_{32}$ está en fila 3, columna 2. Allí aparece $3$, luego $\boxed{a_{32}=3}$.

La entrada $a_{13}$ está en fila 1, columna 3. Allí aparece $2$, luego $\boxed{a_{13}=2}$.


\par\bigskip\noindent\textbf{Solución Ejercicio 41}\par\nobreak
Contamos 2 filas y 5 columnas, así que el orden es $2\times5$. El número de entradas es $2\cdot 5=10$.

La entrada $a_{22}$ está en fila 2, columna 2. Allí aparece $1$, luego $\boxed{a_{22}=1}$.

La entrada $a_{21}$ está en fila 2, columna 1. Allí aparece $3$, luego $\boxed{a_{21}=3}$.


\par\bigskip\noindent\textbf{Solución Ejercicio 43}\par\nobreak
Contamos 5 filas y 3 columnas, así que el orden es $5\times3$. El número de entradas es $5\cdot 3=15$.

La entrada $a_{53}$ está en fila 5, columna 3. Allí aparece $-2$, luego $\boxed{a_{53}=-2}$.

La entrada $a_{22}$ está en fila 2, columna 2. Allí aparece $-3$, luego $\boxed{a_{22}=-3}$.


\par\bigskip\noindent\textbf{Solución Ejercicio 45}\par\nobreak
Contamos 5 filas y 4 columnas, así que el orden es $5\times4$. El número de entradas es $5\cdot 4=20$.

La entrada $a_{23}$ está en fila 2, columna 3. Allí aparece $-2$, luego $\boxed{a_{23}=-2}$.

La entrada $a_{52}$ está en fila 5, columna 2. Allí aparece $-3$, luego $\boxed{a_{52}=-3}$.


\par\bigskip\noindent\textbf{Solución Ejercicio 47}\par\nobreak
Dos matrices del mismo orden son iguales si y sólo si coinciden sus entradas en las mismas posiciones.

En la posición $(1,2)$:
\[
y - 1=2\quad\Longrightarrow\quad y=3.
\]
En la posición $(2,2)$:
\[
x - 2=-3\quad\Longrightarrow\quad x=-1.
\]
Sustituimos estos valores en la matriz del primer miembro:
\[
\begin{pmatrix}-1&2\\-2&-3\end{pmatrix}=\begin{pmatrix}-1&2\\-2&-3\end{pmatrix}.
\]
Como las matrices coinciden entrada por entrada, la solución es $x=-1$, $y=3$.


\par\bigskip\noindent\textbf{Solución Ejercicio 49}\par\nobreak
Dos matrices del mismo orden son iguales si y sólo si coinciden sus entradas en las mismas posiciones.

En la posición $(1,1)$:
\[
x + 1=-3\quad\Longrightarrow\quad x=-4.
\]
En la posición $(3,2)$:
\[
y - 1=1\quad\Longrightarrow\quad y=2.
\]
Sustituimos estos valores en la matriz del primer miembro:
\[
\begin{pmatrix}-3&3\\-2&2\\-1&1\end{pmatrix}=\begin{pmatrix}-3&3\\-2&2\\-1&1\end{pmatrix}.
\]
Como las matrices coinciden entrada por entrada, la solución es $x=-4$, $y=2$.


\par\bigskip\noindent\textbf{Solución Ejercicio 51}\par\nobreak
Dos matrices del mismo orden son iguales si y sólo si coinciden sus entradas en las mismas posiciones.

En la posición $(1,1)$:
\[
x + 2=1\quad\Longrightarrow\quad x=-1.
\]
En la posición $(2,2)$:
\[
y - 1=-3\quad\Longrightarrow\quad y=-2.
\]
Sustituimos estos valores en la matriz del primer miembro:
\[
\begin{pmatrix}1&-3&-1&3\\-3&-3&-1&-1\end{pmatrix}=\begin{pmatrix}1&-3&-1&3\\-3&-3&-1&-1\end{pmatrix}.
\]
Como las matrices coinciden entrada por entrada, la solución es $x=-1$, $y=-2$.


\par\bigskip\noindent\textbf{Solución Ejercicio 53}\par\nobreak
Dos matrices del mismo orden son iguales si y sólo si coinciden sus entradas en las mismas posiciones.

En la posición $(1,2)$:
\[
x + 2=-1\quad\Longrightarrow\quad x=-3.
\]
En la posición $(3,2)$:
\[
y + 1=-3\quad\Longrightarrow\quad y=-4.
\]
Sustituimos estos valores en la matriz del primer miembro:
\[
\begin{pmatrix}-1&-1&3&-2\\-1&-2&-3&0\\3&-3&0&0\end{pmatrix}=\begin{pmatrix}-1&-1&3&-2\\-1&-2&-3&0\\3&-3&0&0\end{pmatrix}.
\]
Como las matrices coinciden entrada por entrada, la solución es $x=-3$, $y=-4$.


\par\bigskip\noindent\textbf{Solución Ejercicio 55}\par\nobreak
Dos matrices del mismo orden son iguales si y sólo si coinciden sus entradas en las mismas posiciones.

En la posición $(2,1)$:
\[
y - 1=-2\quad\Longrightarrow\quad y=-1.
\]
En la posición $(4,1)$:
\[
x - 1=-3\quad\Longrightarrow\quad x=-2.
\]
Sustituimos estos valores en la matriz del primer miembro:
\[
\begin{pmatrix}-3&-3\\-2&-3\\-1&-3\\-3&3\\-2&-3\end{pmatrix}=\begin{pmatrix}-3&-3\\-2&-3\\-1&-3\\-3&3\\-2&-3\end{pmatrix}.
\]
Como las matrices coinciden entrada por entrada, la solución es $x=-2$, $y=-1$.


\par\bigskip\noindent\textbf{Solución Ejercicio 57}\par\nobreak
Dos matrices del mismo orden son iguales si y sólo si coinciden sus entradas en las mismas posiciones.

En la posición $(2,1)$:
\[
y - 2=-1\quad\Longrightarrow\quad y=1.
\]
En la posición $(4,3)$:
\[
x + 1=-2\quad\Longrightarrow\quad x=-3.
\]
Sustituimos estos valores en la matriz del primer miembro:
\[
\begin{pmatrix}0&0&-2&3\\-1&-2&-1&-2\\-1&1&-2&0\\-3&1&-2&-2\end{pmatrix}=\begin{pmatrix}0&0&-2&3\\-1&-2&-1&-2\\-1&1&-2&0\\-3&1&-2&-2\end{pmatrix}.
\]
Como las matrices coinciden entrada por entrada, la solución es $x=-3$, $y=1$.


\par\bigskip\noindent\textbf{Solución Ejercicio 59}\par\nobreak
Primero observamos que tiene 2 filas y 2 columnas. Revisamos las condiciones: todas las entradas fuera de la diagonal son cero; todos los elementos debajo de la diagonal son cero; todos los elementos encima de la diagonal son cero; $A^T=A$.

\textbf{Clasificación:} cuadrada, diagonal, triangular superior, triangular inferior, simétrica.


\par\bigskip\noindent\textbf{Solución Ejercicio 61}\par\nobreak
Primero observamos que tiene 3 filas y 3 columnas. Revisamos las condiciones: hay entradas no nulas fuera de la diagonal; hay elementos no nulos debajo de la diagonal; hay elementos no nulos encima de la diagonal; $A^T=A$.

\textbf{Clasificación:} cuadrada, simétrica.


\par\bigskip\noindent\textbf{Solución Ejercicio 63}\par\nobreak
Primero observamos que tiene 4 filas y 4 columnas. Revisamos las condiciones: hay entradas no nulas fuera de la diagonal; hay elementos no nulos debajo de la diagonal; todos los elementos encima de la diagonal son cero; $A^T\ne A$.

\textbf{Clasificación:} cuadrada, triangular inferior.


\par\bigskip\noindent\textbf{Solución Ejercicio 65}\par\nobreak
Primero observamos que tiene 5 filas y 5 columnas. Revisamos las condiciones: todas las entradas fuera de la diagonal son cero; todos los elementos debajo de la diagonal son cero; todos los elementos encima de la diagonal son cero; $A^T=A$.

\textbf{Clasificación:} cuadrada, diagonal, triangular superior, triangular inferior, simétrica, nula, escalar ($\lambda I$).


\par\bigskip\noindent\textbf{Solución Ejercicio 67}\par\nobreak
Primero observamos que tiene 2 filas y 2 columnas. Revisamos las condiciones: todas las entradas fuera de la diagonal son cero; todos los elementos debajo de la diagonal son cero; todos los elementos encima de la diagonal son cero; $A^T=A$.

\textbf{Clasificación:} cuadrada, diagonal, triangular superior, triangular inferior, simétrica, escalar ($\lambda I$).


\par\bigskip\noindent\textbf{Solución Ejercicio 69}\par\nobreak
Primero observamos que tiene 3 filas y 3 columnas. Revisamos las condiciones: hay entradas no nulas fuera de la diagonal; hay elementos no nulos debajo de la diagonal; hay elementos no nulos encima de la diagonal; $A^T=A$.

\textbf{Clasificación:} cuadrada, simétrica.


\par\bigskip\noindent\textbf{Solución Ejercicio 71}\par\nobreak
Primero multiplicamos cada entrada de $A$ por $2$:
\[
2A=\begin{pmatrix}-4&-2\\4&2\end{pmatrix}.
\]
Después restamos las entradas correspondientes de $B$, sin cambiar el orden:
Fila 1: $(-4)-(-1)=-3\; ;\; (-2)-(-2)=0$.

Fila 2: $(4)-(-2)=6\; ;\; (2)-(1)=1$.

\[
\boxed{2A-B=\begin{pmatrix}-3&0\\6&1\end{pmatrix}}
\]

\par\bigskip\noindent\textbf{Solución Ejercicio 73}\par\nobreak
Primero multiplicamos cada entrada de $A$ por $2$:
\[
2A=\begin{pmatrix}2&2\\-4&2\\2&2\end{pmatrix}.
\]
Después restamos las entradas correspondientes de $B$, sin cambiar el orden:
Fila 1: $(2)-(1)=1\; ;\; (2)-(2)=0$.

Fila 2: $(-4)-(0)=-4\; ;\; (2)-(2)=0$.

Fila 3: $(2)-(1)=1\; ;\; (2)-(-2)=4$.

\[
\boxed{2A-B=\begin{pmatrix}1&0\\-4&0\\1&4\end{pmatrix}}
\]

\par\bigskip\noindent\textbf{Solución Ejercicio 75}\par\nobreak
Primero multiplicamos cada entrada de $A$ por $2$:
\[
2A=\begin{pmatrix}-2&4\\-2&2\\-4&-4\\-4&2\end{pmatrix}.
\]
Después restamos las entradas correspondientes de $B$, sin cambiar el orden:
Fila 1: $(-2)-(1)=-3\; ;\; (4)-(1)=3$.

Fila 2: $(-2)-(1)=-3\; ;\; (2)-(0)=2$.

Fila 3: $(-4)-(2)=-6\; ;\; (-4)-(-2)=-2$.

Fila 4: $(-4)-(-1)=-3\; ;\; (2)-(1)=1$.

\[
\boxed{2A-B=\begin{pmatrix}-3&3\\-3&2\\-6&-2\\-3&1\end{pmatrix}}
\]

\par\bigskip\noindent\textbf{Solución Ejercicio 77}\par\nobreak
Primero multiplicamos cada entrada de $A$ por $2$:
\[
2A=\begin{pmatrix}2&2&4&4\\0&-2&0&4\\-2&2&-2&0\end{pmatrix}.
\]
Después restamos las entradas correspondientes de $B$, sin cambiar el orden:
Fila 1: $(2)-(-2)=4\; ;\; (2)-(2)=0\; ;\; (4)-(0)=4\; ;\; (4)-(1)=3$.

Fila 2: $(0)-(0)=0\; ;\; (-2)-(2)=-4\; ;\; (0)-(2)=-2\; ;\; (4)-(1)=3$.

Fila 3: $(-2)-(1)=-3\; ;\; (2)-(0)=2\; ;\; (-2)-(-1)=-1\; ;\; (0)-(-1)=1$.

\[
\boxed{2A-B=\begin{pmatrix}4&0&4&3\\0&-4&-2&3\\-3&2&-1&1\end{pmatrix}}
\]

\par\bigskip\noindent\textbf{Solución Ejercicio 79}\par\nobreak
Primero multiplicamos cada entrada de $A$ por $2$:
\[
2A=\begin{pmatrix}-2&2\\-4&-4\\2&4\\0&4\\4&-2\end{pmatrix}.
\]
Después restamos las entradas correspondientes de $B$, sin cambiar el orden:
Fila 1: $(-2)-(-1)=-1\; ;\; (2)-(2)=0$.

Fila 2: $(-4)-(2)=-6\; ;\; (-4)-(1)=-5$.

Fila 3: $(2)-(-2)=4\; ;\; (4)-(-1)=5$.

Fila 4: $(0)-(-1)=1\; ;\; (4)-(-1)=5$.

Fila 5: $(4)-(2)=2\; ;\; (-2)-(-2)=0$.

\[
\boxed{2A-B=\begin{pmatrix}-1&0\\-6&-5\\4&5\\1&5\\2&0\end{pmatrix}}
\]

\par\bigskip\noindent\textbf{Solución Ejercicio 81}\par\nobreak
Primero multiplicamos cada entrada de $A$ por $2$:
\[
2A=\begin{pmatrix}2&0&2&-4\\-4&4&2&-4\\4&-2&4&4\\-2&0&2&0\end{pmatrix}.
\]
Después restamos las entradas correspondientes de $B$, sin cambiar el orden:
Fila 1: $(2)-(1)=1\; ;\; (0)-(-2)=2\; ;\; (2)-(1)=1\; ;\; (-4)-(-2)=-2$.

Fila 2: $(-4)-(0)=-4\; ;\; (4)-(1)=3\; ;\; (2)-(0)=2\; ;\; (-4)-(-1)=-3$.

Fila 3: $(4)-(0)=4\; ;\; (-2)-(1)=-3\; ;\; (4)-(-1)=5\; ;\; (4)-(0)=4$.

Fila 4: $(-2)-(2)=-4\; ;\; (0)-(0)=0\; ;\; (2)-(-2)=4\; ;\; (0)-(-1)=1$.

\[
\boxed{2A-B=\begin{pmatrix}1&2&1&-2\\-4&3&2&-3\\4&-3&5&4\\-4&0&4&1\end{pmatrix}}
\]

\par\bigskip\noindent\textbf{Solución Ejercicio 83}\par\nobreak
Primero multiplicamos cada entrada de $A$ por $2$:
\[
2A=\begin{pmatrix}-2&-4&-4&-2&0\\-4&-2&0&-2&-2\\-4&-4&-4&-2&4\end{pmatrix}.
\]
Después restamos las entradas correspondientes de $B$, sin cambiar el orden:
Fila 1: $(-2)-(-1)=-1\; ;\; (-4)-(0)=-4\; ;\; (-4)-(2)=-6\; ;\; (-2)-(0)=-2\; ;\; (0)-(-2)=2$.

Fila 2: $(-4)-(2)=-6\; ;\; (-2)-(-1)=-1\; ;\; (0)-(-1)=1\; ;\; (-2)-(0)=-2\; ;\; (-2)-(2)=-4$.

Fila 3: $(-4)-(0)=-4\; ;\; (-4)-(0)=-4\; ;\; (-4)-(2)=-6\; ;\; (-2)-(-2)=0\; ;\; (4)-(1)=3$.

\[
\boxed{2A-B=\begin{pmatrix}-1&-4&-6&-2&2\\-6&-1&1&-2&-4\\-4&-4&-6&0&3\end{pmatrix}}
\]

\par\bigskip\noindent\textbf{Solución Ejercicio 85}\par\nobreak
Trasponer significa intercambiar filas por columnas: la fila $i$ de $A$ se convierte en la columna $i$ de $A^T$.

\[
A^T=\begin{pmatrix}-1&-1\\2&-1\\0&1\end{pmatrix}.
\]
El orden cambia de $2\times3$ a $3\times2$.

La suma $A+A^T$ \textbf{no está definida}, pues sus sumandos tienen órdenes $2\times3$ y $3\times2$, respectivamente.


\par\bigskip\noindent\textbf{Solución Ejercicio 87}\par\nobreak
Trasponer significa intercambiar filas por columnas: la fila $i$ de $A$ se convierte en la columna $i$ de $A^T$.

\[
A^T=\begin{pmatrix}-1&-2&-2\\-1&0&-1\\-1&-2&-2\end{pmatrix}.
\]
El orden cambia de $3\times3$ a $3\times3$.

Como ambos órdenes coinciden, $A+A^T$ está definida. Para comprobarlo,
\[
A+A^T=\begin{pmatrix}-2&-3&-3\\-3&0&-3\\-3&-3&-4\end{pmatrix}.
\]

\par\bigskip\noindent\textbf{Solución Ejercicio 89}\par\nobreak
Trasponer significa intercambiar filas por columnas: la fila $i$ de $A$ se convierte en la columna $i$ de $A^T$.

\[
A^T=\begin{pmatrix}1&-1\\2&0\\-1&-1\\-1&-2\end{pmatrix}.
\]
El orden cambia de $2\times4$ a $4\times2$.

La suma $A+A^T$ \textbf{no está definida}, pues sus sumandos tienen órdenes $2\times4$ y $4\times2$, respectivamente.


\par\bigskip\noindent\textbf{Solución Ejercicio 91}\par\nobreak
Trasponer significa intercambiar filas por columnas: la fila $i$ de $A$ se convierte en la columna $i$ de $A^T$.

\[
A^T=\begin{pmatrix}0&2&-2&1&-2\\1&1&0&1&-2\end{pmatrix}.
\]
El orden cambia de $5\times2$ a $2\times5$.

La suma $A+A^T$ \textbf{no está definida}, pues sus sumandos tienen órdenes $5\times2$ y $2\times5$, respectivamente.


\par\bigskip\noindent\textbf{Solución Ejercicio 93}\par\nobreak
Trasponer significa intercambiar filas por columnas: la fila $i$ de $A$ se convierte en la columna $i$ de $A^T$.

\[
A^T=\begin{pmatrix}-2&0&-1&2&1\\2&0&1&1&-2\\0&2&1&2&-1\\-1&1&0&-1&0\end{pmatrix}.
\]
El orden cambia de $5\times4$ a $4\times5$.

La suma $A+A^T$ \textbf{no está definida}, pues sus sumandos tienen órdenes $5\times4$ y $4\times5$, respectivamente.


\par\bigskip\noindent\textbf{Solución Ejercicio 95}\par\nobreak
Trasponer significa intercambiar filas por columnas: la fila $i$ de $A$ se convierte en la columna $i$ de $A^T$.

\[
A^T=\begin{pmatrix}-1&2&1&-1&0\\1&-1&2&-1&-1\\-1&0&1&1&-2\\1&-1&-1&1&1\\0&-2&2&-2&2\end{pmatrix}.
\]
El orden cambia de $5\times5$ a $5\times5$.

Como ambos órdenes coinciden, $A+A^T$ está definida. Para comprobarlo,
\[
A+A^T=\begin{pmatrix}-2&3&0&0&0\\3&-2&2&-2&-3\\0&2&2&0&0\\0&-2&0&2&-1\\0&-3&0&-1&4\end{pmatrix}.
\]

\par\bigskip\noindent\textbf{Solución Ejercicio 97}\par\nobreak
El producto $AB$ exige que las 3 columnas de $A$ coincidan con las 3 filas de $B$. Como $3=3$, existe y su orden se forma con las dimensiones exteriores: $\boxed{AB\text{ de orden } 2\times 4}$.

Para $BA$, las 4 columnas de $B$ deben coincidir con las 2 filas de $A$. Como $4\neq2$, $\boxed{BA\text{ no existe}}$.


\par\bigskip\noindent\textbf{Solución Ejercicio 99}\par\nobreak
El producto $AB$ exige que las 3 columnas de $A$ coincidan con las 4 filas de $B$. Como $3\neq4$, $\boxed{AB\text{ no existe}}$.

Para $BA$, las 2 columnas de $B$ deben coincidir con las 4 filas de $A$. Como $2\neq4$, $\boxed{BA\text{ no existe}}$.


\par\bigskip\noindent\textbf{Solución Ejercicio 101}\par\nobreak
El producto $AB$ exige que las 4 columnas de $A$ coincidan con las 4 filas de $B$. Como $4=4$, existe y su orden se forma con las dimensiones exteriores: $\boxed{AB\text{ de orden } 5\times 2}$.

Para $BA$, las 2 columnas de $B$ deben coincidir con las 5 filas de $A$. Como $2\neq5$, $\boxed{BA\text{ no existe}}$.


\par\bigskip\noindent\textbf{Solución Ejercicio 103}\par\nobreak
El producto $AB$ exige que las 2 columnas de $A$ coincidan con las 2 filas de $B$. Como $2=2$, existe y su orden se forma con las dimensiones exteriores: $\boxed{AB\text{ de orden } 4\times 3}$.

Para $BA$, las 3 columnas de $B$ deben coincidir con las 4 filas de $A$. Como $3\neq4$, $\boxed{BA\text{ no existe}}$.


\par\bigskip\noindent\textbf{Solución Ejercicio 105}\par\nobreak
El producto $AB$ exige que las 4 columnas de $A$ coincidan con las 2 filas de $B$. Como $4\neq2$, $\boxed{AB\text{ no existe}}$.

Para $BA$, las 4 columnas de $B$ deben coincidir con las 2 filas de $A$. Como $4\neq2$, $\boxed{BA\text{ no existe}}$.


\par\bigskip\noindent\textbf{Solución Ejercicio 107}\par\nobreak
El producto $AB$ exige que las 5 columnas de $A$ coincidan con las 5 filas de $B$. Como $5=5$, existe y su orden se forma con las dimensiones exteriores: $\boxed{AB\text{ de orden } 1\times 1}$.

Para $BA$, las 1 columnas de $B$ deben coincidir con las 1 filas de $A$. Como $1=1$, existe y su orden es $\boxed{5\times 5}$.


\par\bigskip\noindent\textbf{Solución Ejercicio 109}\par\nobreak
El producto $AB$ exige que las 5 columnas de $A$ coincidan con las 5 filas de $B$. Como $5=5$, existe y su orden se forma con las dimensiones exteriores: $\boxed{AB\text{ de orden } 4\times 4}$.

Para $BA$, las 4 columnas de $B$ deben coincidir con las 4 filas de $A$. Como $4=4$, existe y su orden es $\boxed{5\times 5}$.


\par\bigskip\noindent\textbf{Solución Ejercicio 111}\par\nobreak
El producto $AB$ exige que las 4 columnas de $A$ coincidan con las 5 filas de $B$. Como $4\neq5$, $\boxed{AB\text{ no existe}}$.

Para $BA$, las 2 columnas de $B$ deben coincidir con las 3 filas de $A$. Como $2\neq3$, $\boxed{BA\text{ no existe}}$.


\par\bigskip\noindent\textbf{Solución Ejercicio 113}\par\nobreak
Las dimensiones son $2\times2$ y $2\times2$; como coinciden las dimensiones interiores (2), el producto es $2\times2$.

La entrada $c_{ij}$ se obtiene multiplicando la fila $i$ de $A$ por la columna $j$ de $B$ y sumando los productos:
\[
c_{ij}=\sum_{k=1}^{2} a_{ik}b_{kj}.
\]
\textbf{Fila 1 del producto:}
\[
\begin{aligned}c_{11}&=1\cdot 0+0\cdot (-2)=0\\c_{12}&=1\cdot (-1)+0\cdot (-2)=-1\end{aligned}
\]
\textbf{Fila 2 del producto:}
\[
\begin{aligned}c_{21}&=0\cdot 0+2\cdot (-2)=-4\\c_{22}&=0\cdot (-1)+2\cdot (-2)=-4\end{aligned}
\]
Por tanto,
\[
\boxed{AB=\begin{pmatrix}0&-1\\-4&-4\end{pmatrix}}
\]

\par\bigskip\noindent\textbf{Solución Ejercicio 115}\par\nobreak
Las dimensiones son $3\times2$ y $2\times4$; como coinciden las dimensiones interiores (2), el producto es $3\times4$.

La entrada $c_{ij}$ se obtiene multiplicando la fila $i$ de $A$ por la columna $j$ de $B$ y sumando los productos:
\[
c_{ij}=\sum_{k=1}^{2} a_{ik}b_{kj}.
\]
\textbf{Fila 1 del producto:}
\[
\begin{aligned}c_{11}&=1\cdot (-1)+(-2)\cdot (-1)=1\\c_{12}&=1\cdot 0+(-2)\cdot 2=-4\\c_{13}&=1\cdot (-2)+(-2)\cdot 2=-6\\c_{14}&=1\cdot (-1)+(-2)\cdot 2=-5\end{aligned}
\]
\textbf{Fila 2 del producto:}
\[
\begin{aligned}c_{21}&=0\cdot (-1)+(-2)\cdot (-1)=2\\c_{22}&=0\cdot 0+(-2)\cdot 2=-4\\c_{23}&=0\cdot (-2)+(-2)\cdot 2=-4\\c_{24}&=0\cdot (-1)+(-2)\cdot 2=-4\end{aligned}
\]
\textbf{Fila 3 del producto:}
\[
\begin{aligned}c_{31}&=0\cdot (-1)+(-1)\cdot (-1)=1\\c_{32}&=0\cdot 0+(-1)\cdot 2=-2\\c_{33}&=0\cdot (-2)+(-1)\cdot 2=-2\\c_{34}&=0\cdot (-1)+(-1)\cdot 2=-2\end{aligned}
\]
Por tanto,
\[
\boxed{AB=\begin{pmatrix}1&-4&-6&-5\\2&-4&-4&-4\\1&-2&-2&-2\end{pmatrix}}
\]

\par\bigskip\noindent\textbf{Solución Ejercicio 117}\par\nobreak
Las dimensiones son $4\times2$ y $2\times3$; como coinciden las dimensiones interiores (2), el producto es $4\times3$.

La entrada $c_{ij}$ se obtiene multiplicando la fila $i$ de $A$ por la columna $j$ de $B$ y sumando los productos:
\[
c_{ij}=\sum_{k=1}^{2} a_{ik}b_{kj}.
\]
\textbf{Fila 1 del producto:}
\[
\begin{aligned}c_{11}&=1\cdot 2+1\cdot 1=3\\c_{12}&=1\cdot (-2)+1\cdot (-2)=-4\\c_{13}&=1\cdot (-2)+1\cdot 2=0\end{aligned}
\]
\textbf{Fila 2 del producto:}
\[
\begin{aligned}c_{21}&=2\cdot 2+(-1)\cdot 1=3\\c_{22}&=2\cdot (-2)+(-1)\cdot (-2)=-2\\c_{23}&=2\cdot (-2)+(-1)\cdot 2=-6\end{aligned}
\]
\textbf{Fila 3 del producto:}
\[
\begin{aligned}c_{31}&=2\cdot 2+2\cdot 1=6\\c_{32}&=2\cdot (-2)+2\cdot (-2)=-8\\c_{33}&=2\cdot (-2)+2\cdot 2=0\end{aligned}
\]
\textbf{Fila 4 del producto:}
\[
\begin{aligned}c_{41}&=(-2)\cdot 2+(-1)\cdot 1=-5\\c_{42}&=(-2)\cdot (-2)+(-1)\cdot (-2)=6\\c_{43}&=(-2)\cdot (-2)+(-1)\cdot 2=2\end{aligned}
\]
Por tanto,
\[
\boxed{AB=\begin{pmatrix}3&-4&0\\3&-2&-6\\6&-8&0\\-5&6&2\end{pmatrix}}
\]

\par\bigskip\noindent\textbf{Solución Ejercicio 119}\par\nobreak
Las dimensiones son $4\times3$ y $3\times2$; como coinciden las dimensiones interiores (3), el producto es $4\times2$.

La entrada $c_{ij}$ se obtiene multiplicando la fila $i$ de $A$ por la columna $j$ de $B$ y sumando los productos:
\[
c_{ij}=\sum_{k=1}^{3} a_{ik}b_{kj}.
\]
\textbf{Fila 1 del producto:}
\[
\begin{aligned}c_{11}&=2\cdot 2+2\cdot (-1)+0\cdot 0=2\\c_{12}&=2\cdot (-2)+2\cdot 2+0\cdot 2=0\end{aligned}
\]
\textbf{Fila 2 del producto:}
\[
\begin{aligned}c_{21}&=(-1)\cdot 2+0\cdot (-1)+1\cdot 0=-2\\c_{22}&=(-1)\cdot (-2)+0\cdot 2+1\cdot 2=4\end{aligned}
\]
\textbf{Fila 3 del producto:}
\[
\begin{aligned}c_{31}&=(-2)\cdot 2+2\cdot (-1)+0\cdot 0=-6\\c_{32}&=(-2)\cdot (-2)+2\cdot 2+0\cdot 2=8\end{aligned}
\]
\textbf{Fila 4 del producto:}
\[
\begin{aligned}c_{41}&=(-1)\cdot 2+1\cdot (-1)+(-2)\cdot 0=-3\\c_{42}&=(-1)\cdot (-2)+1\cdot 2+(-2)\cdot 2=0\end{aligned}
\]
Por tanto,
\[
\boxed{AB=\begin{pmatrix}2&0\\-2&4\\-6&8\\-3&0\end{pmatrix}}
\]

\par\bigskip\noindent\textbf{Solución Ejercicio 121}\par\nobreak
Las dimensiones son $3\times5$ y $5\times2$; como coinciden las dimensiones interiores (5), el producto es $3\times2$.

La entrada $c_{ij}$ se obtiene multiplicando la fila $i$ de $A$ por la columna $j$ de $B$ y sumando los productos:
\[
c_{ij}=\sum_{k=1}^{5} a_{ik}b_{kj}.
\]
\textbf{Fila 1 del producto:}
\[
\begin{aligned}c_{11}&=1\cdot 2+(-2)\cdot 0+0\cdot 0+0\cdot 2+1\cdot (-2)=0\\c_{12}&=1\cdot 0+(-2)\cdot (-1)+0\cdot 2+0\cdot 1+1\cdot 0=2\end{aligned}
\]
\textbf{Fila 2 del producto:}
\[
\begin{aligned}c_{21}&=(-1)\cdot 2+(-1)\cdot 0+(-2)\cdot 0+0\cdot 2+(-2)\cdot (-2)=2\\c_{22}&=(-1)\cdot 0+(-1)\cdot (-1)+(-2)\cdot 2+0\cdot 1+(-2)\cdot 0=-3\end{aligned}
\]
\textbf{Fila 3 del producto:}
\[
\begin{aligned}c_{31}&=2\cdot 2+0\cdot 0+(-2)\cdot 0+2\cdot 2+1\cdot (-2)=6\\c_{32}&=2\cdot 0+0\cdot (-1)+(-2)\cdot 2+2\cdot 1+1\cdot 0=-2\end{aligned}
\]
Por tanto,
\[
\boxed{AB=\begin{pmatrix}0&2\\2&-3\\6&-2\end{pmatrix}}
\]

\par\bigskip\noindent\textbf{Solución Ejercicio 123}\par\nobreak
Las dimensiones son $5\times3$ y $3\times4$; como coinciden las dimensiones interiores (3), el producto es $5\times4$.

La entrada $c_{ij}$ se obtiene multiplicando la fila $i$ de $A$ por la columna $j$ de $B$ y sumando los productos:
\[
c_{ij}=\sum_{k=1}^{3} a_{ik}b_{kj}.
\]
\textbf{Fila 1 del producto:}
\[
\begin{aligned}c_{11}&=2\cdot 1+2\cdot 0+2\cdot (-1)=0\\c_{12}&=2\cdot (-1)+2\cdot (-2)+2\cdot 1=-4\\c_{13}&=2\cdot 1+2\cdot (-1)+2\cdot (-2)=-4\\c_{14}&=2\cdot 0+2\cdot (-2)+2\cdot 1=-2\end{aligned}
\]
\textbf{Fila 2 del producto:}
\[
\begin{aligned}c_{21}&=1\cdot 1+(-2)\cdot 0+2\cdot (-1)=-1\\c_{22}&=1\cdot (-1)+(-2)\cdot (-2)+2\cdot 1=5\\c_{23}&=1\cdot 1+(-2)\cdot (-1)+2\cdot (-2)=-1\\c_{24}&=1\cdot 0+(-2)\cdot (-2)+2\cdot 1=6\end{aligned}
\]
\textbf{Fila 3 del producto:}
\[
\begin{aligned}c_{31}&=0\cdot 1+2\cdot 0+(-1)\cdot (-1)=1\\c_{32}&=0\cdot (-1)+2\cdot (-2)+(-1)\cdot 1=-5\\c_{33}&=0\cdot 1+2\cdot (-1)+(-1)\cdot (-2)=0\\c_{34}&=0\cdot 0+2\cdot (-2)+(-1)\cdot 1=-5\end{aligned}
\]
\textbf{Fila 4 del producto:}
\[
\begin{aligned}c_{41}&=2\cdot 1+(-2)\cdot 0+(-1)\cdot (-1)=3\\c_{42}&=2\cdot (-1)+(-2)\cdot (-2)+(-1)\cdot 1=1\\c_{43}&=2\cdot 1+(-2)\cdot (-1)+(-1)\cdot (-2)=6\\c_{44}&=2\cdot 0+(-2)\cdot (-2)+(-1)\cdot 1=3\end{aligned}
\]
\textbf{Fila 5 del producto:}
\[
\begin{aligned}c_{51}&=2\cdot 1+1\cdot 0+2\cdot (-1)=0\\c_{52}&=2\cdot (-1)+1\cdot (-2)+2\cdot 1=-2\\c_{53}&=2\cdot 1+1\cdot (-1)+2\cdot (-2)=-3\\c_{54}&=2\cdot 0+1\cdot (-2)+2\cdot 1=0\end{aligned}
\]
Por tanto,
\[
\boxed{AB=\begin{pmatrix}0&-4&-4&-2\\-1&5&-1&6\\1&-5&0&-5\\3&1&6&3\\0&-2&-3&0\end{pmatrix}}
\]

\par\bigskip\noindent\textbf{Solución Ejercicio 125}\par\nobreak
Las dimensiones son $5\times5$ y $5\times2$; como coinciden las dimensiones interiores (5), el producto es $5\times2$.

La entrada $c_{ij}$ se obtiene multiplicando la fila $i$ de $A$ por la columna $j$ de $B$ y sumando los productos:
\[
c_{ij}=\sum_{k=1}^{5} a_{ik}b_{kj}.
\]
\textbf{Fila 1 del producto:}
\[
\begin{aligned}c_{11}&=(-1)\cdot 0+2\cdot (-1)+0\cdot (-1)+(-2)\cdot (-2)+(-1)\cdot 0=2\\c_{12}&=(-1)\cdot 1+2\cdot (-1)+0\cdot (-2)+(-2)\cdot (-2)+(-1)\cdot 2=-1\end{aligned}
\]
\textbf{Fila 2 del producto:}
\[
\begin{aligned}c_{21}&=0\cdot 0+2\cdot (-1)+2\cdot (-1)+(-1)\cdot (-2)+2\cdot 0=-2\\c_{22}&=0\cdot 1+2\cdot (-1)+2\cdot (-2)+(-1)\cdot (-2)+2\cdot 2=0\end{aligned}
\]
\textbf{Fila 3 del producto:}
\[
\begin{aligned}c_{31}&=(-2)\cdot 0+(-2)\cdot (-1)+(-2)\cdot (-1)+0\cdot (-2)+(-2)\cdot 0=4\\c_{32}&=(-2)\cdot 1+(-2)\cdot (-1)+(-2)\cdot (-2)+0\cdot (-2)+(-2)\cdot 2=0\end{aligned}
\]
\textbf{Fila 4 del producto:}
\[
\begin{aligned}c_{41}&=1\cdot 0+0\cdot (-1)+(-2)\cdot (-1)+0\cdot (-2)+0\cdot 0=2\\c_{42}&=1\cdot 1+0\cdot (-1)+(-2)\cdot (-2)+0\cdot (-2)+0\cdot 2=5\end{aligned}
\]
\textbf{Fila 5 del producto:}
\[
\begin{aligned}c_{51}&=(-1)\cdot 0+(-2)\cdot (-1)+(-1)\cdot (-1)+1\cdot (-2)+2\cdot 0=1\\c_{52}&=(-1)\cdot 1+(-2)\cdot (-1)+(-1)\cdot (-2)+1\cdot (-2)+2\cdot 2=5\end{aligned}
\]
Por tanto,
\[
\boxed{AB=\begin{pmatrix}2&-1\\-2&0\\4&0\\2&5\\1&5\end{pmatrix}}
\]

\par\bigskip\noindent\textbf{Solución Ejercicio 127}\par\nobreak
Comenzamos con la matriz aumentada $[A\mid I]$ y aplicamos operaciones elementales; cada operación actúa sobre toda la fila.

\[
\left[\begin{array}{rr|rr}1&0&1&0\\2&1&0&1\end{array}\right]
\]
Operación $R_{2}\leftarrow R_{2}-2R_{1}$:
\[
\left[\begin{array}{rr|rr}1&0&1&0\\0&1&-2&1\end{array}\right]
\]
Como la parte izquierda es $I$, la parte derecha es la inversa:
\[
\boxed{A^{-1}=\begin{pmatrix}1&0\\-2&1\end{pmatrix}}
\]
Como comprobación, multiplicamos y obtenemos
\[
AA^{-1}=\begin{pmatrix}1&0\\2&1\end{pmatrix}\begin{pmatrix}1&0\\-2&1\end{pmatrix}=\begin{pmatrix}1&0\\0&1\end{pmatrix}=I_{2}.
\]

\par\bigskip\noindent\textbf{Solución Ejercicio 129}\par\nobreak
Comenzamos con la matriz aumentada $[A\mid I]$ y aplicamos operaciones elementales; cada operación actúa sobre toda la fila.

\[
\left[\begin{array}{rr|rr}-3&2&1&0\\-2&1&0&1\end{array}\right]
\]
Operación $R_{1}\leftarrow - \frac{1}{3}R_{1}$:
\[
\left[\begin{array}{rr|rr}1&- \frac{2}{3}&- \frac{1}{3}&0\\-2&1&0&1\end{array}\right]
\]
Operación $R_{2}\leftarrow R_{2}+2R_{1}$:
\[
\left[\begin{array}{rr|rr}1&- \frac{2}{3}&- \frac{1}{3}&0\\0&- \frac{1}{3}&- \frac{2}{3}&1\end{array}\right]
\]
Operación $R_{2}\leftarrow -3R_{2}$:
\[
\left[\begin{array}{rr|rr}1&- \frac{2}{3}&- \frac{1}{3}&0\\0&1&2&-3\end{array}\right]
\]
Operación $R_{1}\leftarrow R_{1}+\frac{2}{3}R_{2}$:
\[
\left[\begin{array}{rr|rr}1&0&1&-2\\0&1&2&-3\end{array}\right]
\]
Como la parte izquierda es $I$, la parte derecha es la inversa:
\[
\boxed{A^{-1}=\begin{pmatrix}1&-2\\2&-3\end{pmatrix}}
\]
Como comprobación, multiplicamos y obtenemos
\[
AA^{-1}=\begin{pmatrix}-3&2\\-2&1\end{pmatrix}\begin{pmatrix}1&-2\\2&-3\end{pmatrix}=\begin{pmatrix}1&0\\0&1\end{pmatrix}=I_{2}.
\]

\par\bigskip\noindent\textbf{Solución Ejercicio 131}\par\nobreak
Comenzamos con la matriz aumentada $[A\mid I]$ y aplicamos operaciones elementales; cada operación actúa sobre toda la fila.

\[
\left[\begin{array}{rrr|rrr}-3&2&0&1&0&0\\-2&1&0&0&1&0\\2&0&1&0&0&1\end{array}\right]
\]
Operación $R_{1}\leftarrow - \frac{1}{3}R_{1}$:
\[
\left[\begin{array}{rrr|rrr}1&- \frac{2}{3}&0&- \frac{1}{3}&0&0\\-2&1&0&0&1&0\\2&0&1&0&0&1\end{array}\right]
\]
Operación $R_{2}\leftarrow R_{2}+2R_{1}$:
\[
\left[\begin{array}{rrr|rrr}1&- \frac{2}{3}&0&- \frac{1}{3}&0&0\\0&- \frac{1}{3}&0&- \frac{2}{3}&1&0\\2&0&1&0&0&1\end{array}\right]
\]
Operación $R_{3}\leftarrow R_{3}-2R_{1}$:
\[
\left[\begin{array}{rrr|rrr}1&- \frac{2}{3}&0&- \frac{1}{3}&0&0\\0&- \frac{1}{3}&0&- \frac{2}{3}&1&0\\0&\frac{4}{3}&1&\frac{2}{3}&0&1\end{array}\right]
\]
Operación $R_{2}\leftarrow -3R_{2}$:
\[
\left[\begin{array}{rrr|rrr}1&- \frac{2}{3}&0&- \frac{1}{3}&0&0\\0&1&0&2&-3&0\\0&\frac{4}{3}&1&\frac{2}{3}&0&1\end{array}\right]
\]
Operación $R_{1}\leftarrow R_{1}+\frac{2}{3}R_{2}$:
\[
\left[\begin{array}{rrr|rrr}1&0&0&1&-2&0\\0&1&0&2&-3&0\\0&\frac{4}{3}&1&\frac{2}{3}&0&1\end{array}\right]
\]
Operación $R_{3}\leftarrow R_{3}-\frac{4}{3}R_{2}$:
\[
\left[\begin{array}{rrr|rrr}1&0&0&1&-2&0\\0&1&0&2&-3&0\\0&0&1&-2&4&1\end{array}\right]
\]
Como la parte izquierda es $I$, la parte derecha es la inversa:
\[
\boxed{A^{-1}=\begin{pmatrix}1&-2&0\\2&-3&0\\-2&4&1\end{pmatrix}}
\]
Como comprobación, multiplicamos y obtenemos
\[
AA^{-1}=\begin{pmatrix}-3&2&0\\-2&1&0\\2&0&1\end{pmatrix}\begin{pmatrix}1&-2&0\\2&-3&0\\-2&4&1\end{pmatrix}=\begin{pmatrix}1&0&0\\0&1&0\\0&0&1\end{pmatrix}=I_{3}.
\]

\par\bigskip\noindent\textbf{Solución Ejercicio 133}\par\nobreak
Comenzamos con la matriz aumentada $[A\mid I]$ y aplicamos operaciones elementales; cada operación actúa sobre toda la fila.

\[
\left[\begin{array}{rrr|rrr}1&0&0&1&0&0\\0&1&-1&0&1&0\\0&0&1&0&0&1\end{array}\right]
\]
Operación $R_{2}\leftarrow R_{2}+1R_{3}$:
\[
\left[\begin{array}{rrr|rrr}1&0&0&1&0&0\\0&1&0&0&1&1\\0&0&1&0&0&1\end{array}\right]
\]
Como la parte izquierda es $I$, la parte derecha es la inversa:
\[
\boxed{A^{-1}=\begin{pmatrix}1&0&0\\0&1&1\\0&0&1\end{pmatrix}}
\]
Como comprobación, multiplicamos y obtenemos
\[
AA^{-1}=\begin{pmatrix}1&0&0\\0&1&-1\\0&0&1\end{pmatrix}\begin{pmatrix}1&0&0\\0&1&1\\0&0&1\end{pmatrix}=\begin{pmatrix}1&0&0\\0&1&0\\0&0&1\end{pmatrix}=I_{3}.
\]

\par\bigskip\noindent\textbf{Solución Ejercicio 135}\par\nobreak
Comenzamos con la matriz aumentada $[A\mid I]$ y aplicamos operaciones elementales; cada operación actúa sobre toda la fila.

\[
\left[\begin{array}{rrrr|rrrr}1&2&0&0&1&0&0&0\\0&1&0&0&0&1&0&0\\0&0&1&-1&0&0&1&0\\0&0&0&1&0&0&0&1\end{array}\right]
\]
Operación $R_{1}\leftarrow R_{1}-2R_{2}$:
\[
\left[\begin{array}{rrrr|rrrr}1&0&0&0&1&-2&0&0\\0&1&0&0&0&1&0&0\\0&0&1&-1&0&0&1&0\\0&0&0&1&0&0&0&1\end{array}\right]
\]
Operación $R_{3}\leftarrow R_{3}+1R_{4}$:
\[
\left[\begin{array}{rrrr|rrrr}1&0&0&0&1&-2&0&0\\0&1&0&0&0&1&0&0\\0&0&1&0&0&0&1&1\\0&0&0&1&0&0&0&1\end{array}\right]
\]
Como la parte izquierda es $I$, la parte derecha es la inversa:
\[
\boxed{A^{-1}=\begin{pmatrix}1&-2&0&0\\0&1&0&0\\0&0&1&1\\0&0&0&1\end{pmatrix}}
\]
Como comprobación, multiplicamos y obtenemos
\[
AA^{-1}=\begin{pmatrix}1&2&0&0\\0&1&0&0\\0&0&1&-1\\0&0&0&1\end{pmatrix}\begin{pmatrix}1&-2&0&0\\0&1&0&0\\0&0&1&1\\0&0&0&1\end{pmatrix}=\begin{pmatrix}1&0&0&0\\0&1&0&0\\0&0&1&0\\0&0&0&1\end{pmatrix}=I_{4}.
\]

\par\bigskip\noindent\textbf{Solución Ejercicio 137}\par\nobreak
Comenzamos con la matriz aumentada $[A\mid I]$ y aplicamos operaciones elementales; cada operación actúa sobre toda la fila.

\[
\left[\begin{array}{rrrr|rrrr}1&0&0&1&1&0&0&0\\0&1&0&0&0&1&0&0\\0&0&1&0&0&0&1&0\\0&0&-1&1&0&0&0&1\end{array}\right]
\]
Operación $R_{4}\leftarrow R_{4}+1R_{3}$:
\[
\left[\begin{array}{rrrr|rrrr}1&0&0&1&1&0&0&0\\0&1&0&0&0&1&0&0\\0&0&1&0&0&0&1&0\\0&0&0&1&0&0&1&1\end{array}\right]
\]
Operación $R_{1}\leftarrow R_{1}-1R_{4}$:
\[
\left[\begin{array}{rrrr|rrrr}1&0&0&0&1&0&-1&-1\\0&1&0&0&0&1&0&0\\0&0&1&0&0&0&1&0\\0&0&0&1&0&0&1&1\end{array}\right]
\]
Como la parte izquierda es $I$, la parte derecha es la inversa:
\[
\boxed{A^{-1}=\begin{pmatrix}1&0&-1&-1\\0&1&0&0\\0&0&1&0\\0&0&1&1\end{pmatrix}}
\]
Como comprobación, multiplicamos y obtenemos
\[
AA^{-1}=\begin{pmatrix}1&0&0&1\\0&1&0&0\\0&0&1&0\\0&0&-1&1\end{pmatrix}\begin{pmatrix}1&0&-1&-1\\0&1&0&0\\0&0&1&0\\0&0&1&1\end{pmatrix}=\begin{pmatrix}1&0&0&0\\0&1&0&0\\0&0&1&0\\0&0&0&1\end{pmatrix}=I_{4}.
\]

\par\bigskip\noindent\textbf{Solución Ejercicio 139}\par\nobreak
Usamos $\det\begin{psmallmatrix}a&b\\c&d\end{psmallmatrix}=ad-bc$:
\[
\begin{vmatrix}-4&-4\\4&-4\end{vmatrix}=(-4)\cdot (-4)-(-4)\cdot 4=16-(-16)=32.
\]

\par\bigskip\noindent\textbf{Solución Ejercicio 141}\par\nobreak
Usamos $\det\begin{psmallmatrix}a&b\\c&d\end{psmallmatrix}=ad-bc$:
\[
\begin{vmatrix}3&2\\-1&-2\end{vmatrix}=3\cdot (-2)-2\cdot (-1)=-6-(-2)=-4.
\]

\par\bigskip\noindent\textbf{Solución Ejercicio 143}\par\nobreak
Usamos $\det\begin{psmallmatrix}a&b\\c&d\end{psmallmatrix}=ad-bc$:
\[
\begin{vmatrix}2&-1\\0&2\end{vmatrix}=2\cdot 2-(-1)\cdot 0=4-(0)=4.
\]

\par\bigskip\noindent\textbf{Solución Ejercicio 145}\par\nobreak
Usamos $\det\begin{psmallmatrix}a&b\\c&d\end{psmallmatrix}=ad-bc$:
\[
\begin{vmatrix}1&-1\\-3&-2\end{vmatrix}=1\cdot (-2)-(-1)\cdot (-3)=-2-(3)=-5.
\]

\par\bigskip\noindent\textbf{Solución Ejercicio 147}\par\nobreak
Usamos $\det\begin{psmallmatrix}a&b\\c&d\end{psmallmatrix}=ad-bc$:
\[
\begin{vmatrix}-1&-3\\-1&2\end{vmatrix}=(-1)\cdot 2-(-3)\cdot (-1)=-2-(3)=-5.
\]

\par\bigskip\noindent\textbf{Solución Ejercicio 149}\par\nobreak
Usamos $\det\begin{psmallmatrix}a&b\\c&d\end{psmallmatrix}=ad-bc$:
\[
\begin{vmatrix}2&4\\3&-2\end{vmatrix}=2\cdot (-2)-4\cdot 3=-4-(12)=-16.
\]

\par\bigskip\noindent\textbf{Solución Ejercicio 151}\par\nobreak
Consideremos $A=\begin{pmatrix}2&-2&0\\3&3&0\\2&1&3\end{pmatrix}$.

Elegimos la columna 3, que contiene dos ceros, para reducir el trabajo. El patrón de signos es $C_{ij}=(-1)^{i+j}M_{ij}$, donde $M_{ij}$ es el determinante que queda al borrar la fila $i$ y la columna $j$.

Para la entrada $a_{33}=3$, borramos su fila y su columna:
\[
M_{33}=\begin{vmatrix}2&-2\\3&3\end{vmatrix}.
\]
\[
M_{33}=2\cdot 3-(-2)\cdot 3=12.
\]
\[
C_{33}=(-1)^{3+3}(12)=12.
\]
Expandimos: \[
\boxed{\det(A)=3\cdot 12=36}
\]

\par\bigskip\noindent\textbf{Solución Ejercicio 153}\par\nobreak
Consideremos $A=\begin{pmatrix}3&3&-1\\3&0&1\\-2&-1&-2\end{pmatrix}$.

Elegimos la fila 2, que contiene un cero, para reducir el trabajo. El patrón de signos es $C_{ij}=(-1)^{i+j}M_{ij}$, donde $M_{ij}$ es el determinante que queda al borrar la fila $i$ y la columna $j$.

Para la entrada $a_{21}=3$, borramos su fila y su columna:
\[
M_{21}=\begin{vmatrix}3&-1\\-1&-2\end{vmatrix}.
\]
\[
M_{21}=3\cdot (-2)-(-1)\cdot (-1)=-7.
\]
\[
C_{21}=(-1)^{2+1}(-7)=7.
\]
Para la entrada $a_{23}=1$, borramos su fila y su columna:
\[
M_{23}=\begin{vmatrix}3&3\\-2&-1\end{vmatrix}.
\]
\[
M_{23}=3\cdot (-1)-3\cdot (-2)=3.
\]
\[
C_{23}=(-1)^{2+3}(3)=-3.
\]
Expandimos: \[
\boxed{\det(A)=3\cdot 7 + 1\cdot (-3)=18}
\]

\par\bigskip\noindent\textbf{Solución Ejercicio 155}\par\nobreak
Consideremos $A=\begin{pmatrix}1&-1&2\\0&-2&1\\1&2&3\end{pmatrix}$.

Elegimos la fila 2, que contiene un cero, para reducir el trabajo. El patrón de signos es $C_{ij}=(-1)^{i+j}M_{ij}$, donde $M_{ij}$ es el determinante que queda al borrar la fila $i$ y la columna $j$.

Para la entrada $a_{22}=-2$, borramos su fila y su columna:
\[
M_{22}=\begin{vmatrix}1&2\\1&3\end{vmatrix}.
\]
\[
M_{22}=1\cdot 3-2\cdot 1=1.
\]
\[
C_{22}=(-1)^{2+2}(1)=1.
\]
Para la entrada $a_{23}=1$, borramos su fila y su columna:
\[
M_{23}=\begin{vmatrix}1&-1\\1&2\end{vmatrix}.
\]
\[
M_{23}=1\cdot 2-(-1)\cdot 1=3.
\]
\[
C_{23}=(-1)^{2+3}(3)=-3.
\]
Expandimos: \[
\boxed{\det(A)=(-2)\cdot 1 + 1\cdot (-3)=-5}
\]

\par\bigskip\noindent\textbf{Solución Ejercicio 157}\par\nobreak
Consideremos $A=\begin{pmatrix}-1&3&0\\0&-1&2\\2&-1&-2\end{pmatrix}$.

Elegimos la fila 1, que contiene un cero, para reducir el trabajo. El patrón de signos es $C_{ij}=(-1)^{i+j}M_{ij}$, donde $M_{ij}$ es el determinante que queda al borrar la fila $i$ y la columna $j$.

Para la entrada $a_{11}=-1$, borramos su fila y su columna:
\[
M_{11}=\begin{vmatrix}-1&2\\-1&-2\end{vmatrix}.
\]
\[
M_{11}=(-1)\cdot (-2)-2\cdot (-1)=4.
\]
\[
C_{11}=(-1)^{1+1}(4)=4.
\]
Para la entrada $a_{12}=3$, borramos su fila y su columna:
\[
M_{12}=\begin{vmatrix}0&2\\2&-2\end{vmatrix}.
\]
\[
M_{12}=0\cdot (-2)-2\cdot 2=-4.
\]
\[
C_{12}=(-1)^{1+2}(-4)=4.
\]
Expandimos: \[
\boxed{\det(A)=(-1)\cdot 4 + 3\cdot 4=8}
\]

\par\bigskip\noindent\textbf{Solución Ejercicio 159}\par\nobreak
Consideremos $A=\begin{pmatrix}2&2&-3\\0&2&3\\0&3&-1\end{pmatrix}$.

Elegimos la columna 1, que contiene dos ceros, para reducir el trabajo. El patrón de signos es $C_{ij}=(-1)^{i+j}M_{ij}$, donde $M_{ij}$ es el determinante que queda al borrar la fila $i$ y la columna $j$.

Para la entrada $a_{11}=2$, borramos su fila y su columna:
\[
M_{11}=\begin{vmatrix}2&3\\3&-1\end{vmatrix}.
\]
\[
M_{11}=2\cdot (-1)-3\cdot 3=-11.
\]
\[
C_{11}=(-1)^{1+1}(-11)=-11.
\]
Expandimos: \[
\boxed{\det(A)=2\cdot (-11)=-22}
\]

\par\bigskip\noindent\textbf{Solución Ejercicio 161}\par\nobreak
Consideremos $A=\begin{pmatrix}-3&-2&1\\2&2&-3\\1&-3&0\end{pmatrix}$.

Elegimos la fila 3, que contiene un cero, para reducir el trabajo. El patrón de signos es $C_{ij}=(-1)^{i+j}M_{ij}$, donde $M_{ij}$ es el determinante que queda al borrar la fila $i$ y la columna $j$.

Para la entrada $a_{31}=1$, borramos su fila y su columna:
\[
M_{31}=\begin{vmatrix}-2&1\\2&-3\end{vmatrix}.
\]
\[
M_{31}=(-2)\cdot (-3)-1\cdot 2=4.
\]
\[
C_{31}=(-1)^{3+1}(4)=4.
\]
Para la entrada $a_{32}=-3$, borramos su fila y su columna:
\[
M_{32}=\begin{vmatrix}-3&1\\2&-3\end{vmatrix}.
\]
\[
M_{32}=(-3)\cdot (-3)-1\cdot 2=7.
\]
\[
C_{32}=(-1)^{3+2}(7)=-7.
\]
Expandimos: \[
\boxed{\det(A)=1\cdot 4 + (-3)\cdot (-7)=25}
\]

\par\bigskip\noindent\textbf{Solución Ejercicio 163}\par\nobreak
Es triangular; las entradas bajo (o sobre) la diagonal principal son cero. El determinante es el producto de las entradas de esa diagonal:
\[
\det(A)=2\cdot (-1)\cdot (-2)\cdot 3=12.
\]

\par\bigskip\noindent\textbf{Solución Ejercicio 165}\par\nobreak
Consideremos $A=\begin{pmatrix}0&0&1&0\\0&-1&0&0\\-2&1&0&0\\0&0&0&2\end{pmatrix}$.

Elegimos la fila 1, que contiene tres ceros, para reducir el trabajo. El patrón de signos es $C_{ij}=(-1)^{i+j}M_{ij}$, donde $M_{ij}$ es el determinante que queda al borrar la fila $i$ y la columna $j$.

Para la entrada $a_{13}=1$, borramos su fila y su columna:
\[
M_{13}=\begin{vmatrix}0&-1&0\\-2&1&0\\0&0&2\end{vmatrix}.
\]
\[
\begin{aligned}M_{13}&=(-1)\cdot (-1)^{1+2}\cdot\bigl((-2)\cdot 2-0\cdot 0\bigr)\\&=-4\end{aligned}
\]
\[
C_{13}=(-1)^{1+3}(-4)=-4.
\]
Expandimos: \[
\boxed{\det(A)=1\cdot (-4)=-4}
\]

\par\bigskip\noindent\textbf{Solución Ejercicio 167}\par\nobreak
Consideremos $A=\begin{pmatrix}1&0&0&0\\-1&-1&0&-1\\0&0&1&0\\0&0&0&3\end{pmatrix}$.

Elegimos la fila 1, que contiene tres ceros, para reducir el trabajo. El patrón de signos es $C_{ij}=(-1)^{i+j}M_{ij}$, donde $M_{ij}$ es el determinante que queda al borrar la fila $i$ y la columna $j$.

Para la entrada $a_{11}=1$, borramos su fila y su columna:
\[
M_{11}=\begin{vmatrix}-1&0&-1\\0&1&0\\0&0&3\end{vmatrix}.
\]
\[
\begin{aligned}M_{11}&=1\cdot (-1)^{2+2}\cdot\bigl((-1)\cdot 3-(-1)\cdot 0\bigr)\\&=-3\end{aligned}
\]
\[
C_{11}=(-1)^{1+1}(-3)=-3.
\]
Expandimos: \[
\boxed{\det(A)=1\cdot (-3)=-3}
\]

\par\bigskip\noindent\textbf{Solución Ejercicio 169}\par\nobreak
Es triangular; las entradas bajo (o sobre) la diagonal principal son cero. El determinante es el producto de las entradas de esa diagonal:
\[
\det(A)=(-2)\cdot 2\cdot 3\cdot (-1)\cdot 2=24.
\]

\par\bigskip\noindent\textbf{Solución Ejercicio 171}\par\nobreak
Es triangular; las entradas bajo (o sobre) la diagonal principal son cero. El determinante es el producto de las entradas de esa diagonal:
\[
\det(A)=(-1)\cdot 2\cdot 1\cdot 2\cdot 1=-4.
\]

\par\bigskip\noindent\textbf{Solución Ejercicio 173}\par\nobreak
Es triangular; las entradas bajo (o sobre) la diagonal principal son cero. El determinante es el producto de las entradas de esa diagonal:
\[
\det(A)=(-2)\cdot (-1)\cdot 1\cdot 1\cdot 3=6.
\]

\par\bigskip\noindent\textbf{Solución Ejercicio 175}\par\nobreak
Consideremos $A=\begin{pmatrix}-2&2&2\\1&-2&-2\\-2&2&-1\end{pmatrix}$.

El elemento indicado ocupa la fila 3 y la columna 3. Primero \textbf{eliminamos esa fila y esa columna}.

\[
M_{33}=\begin{vmatrix}-2&2\\1&-2\end{vmatrix}.
\]
\[
M_{33}=(-2)\cdot (-2)-2\cdot 1=2.
\]
El \textbf{menor} es el número anterior. El \textbf{cofactor} incorpora además el signo $(-1)^{i+j}$:
\[
\boxed{C_{33}=(-1)^{3+3}M_{33}=1\cdot 2=2}.
\]

\par\bigskip\noindent\textbf{Solución Ejercicio 177}\par\nobreak
Consideremos $A=\begin{pmatrix}1&-1&-1\\0&-2&-2\\0&1&-1\end{pmatrix}$.

El elemento indicado ocupa la fila 3 y la columna 3. Primero \textbf{eliminamos esa fila y esa columna}.

\[
M_{33}=\begin{vmatrix}1&-1\\0&-2\end{vmatrix}.
\]
\[
M_{33}=1\cdot (-2)-(-1)\cdot 0=-2.
\]
El \textbf{menor} es el número anterior. El \textbf{cofactor} incorpora además el signo $(-1)^{i+j}$:
\[
\boxed{C_{33}=(-1)^{3+3}M_{33}=1\cdot (-2)=-2}.
\]

\par\bigskip\noindent\textbf{Solución Ejercicio 179}\par\nobreak
Consideremos $A=\begin{pmatrix}2&0&-2\\1&-1&-1\\-2&0&-1\end{pmatrix}$.

El elemento indicado ocupa la fila 3 y la columna 2. Primero \textbf{eliminamos esa fila y esa columna}.

\[
M_{32}=\begin{vmatrix}2&-2\\1&-1\end{vmatrix}.
\]
\[
M_{32}=2\cdot (-1)-(-2)\cdot 1=0.
\]
El \textbf{menor} es el número anterior. El \textbf{cofactor} incorpora además el signo $(-1)^{i+j}$:
\[
\boxed{C_{32}=(-1)^{3+2}M_{32}=(-1)\cdot 0=0}.
\]

\par\bigskip\noindent\textbf{Solución Ejercicio 181}\par\nobreak
Consideremos $A=\begin{pmatrix}-2&2&-2\\0&2&2\\2&2&0\end{pmatrix}$.

El elemento indicado ocupa la fila 3 y la columna 2. Primero \textbf{eliminamos esa fila y esa columna}.

\[
M_{32}=\begin{vmatrix}-2&-2\\0&2\end{vmatrix}.
\]
\[
M_{32}=(-2)\cdot 2-(-2)\cdot 0=-4.
\]
El \textbf{menor} es el número anterior. El \textbf{cofactor} incorpora además el signo $(-1)^{i+j}$:
\[
\boxed{C_{32}=(-1)^{3+2}M_{32}=(-1)\cdot (-4)=4}.
\]

\par\bigskip\noindent\textbf{Solución Ejercicio 183}\par\nobreak
Consideremos $A=\begin{pmatrix}1&0&1&-1\\-2&1&-2&-2\\-2&-1&-1&2\\-1&0&2&-2\end{pmatrix}$.

El elemento indicado ocupa la fila 4 y la columna 2. Primero \textbf{eliminamos esa fila y esa columna}.

\[
M_{42}=\begin{vmatrix}1&1&-1\\-2&-2&-2\\-2&-1&2\end{vmatrix}.
\]
Calculamos este determinante de orden tres por cofactores:
\[
\begin{aligned}M_{42}&=1\cdot (-1)^{1+1}\cdot\bigl((-2)\cdot 2-(-2)\cdot (-1)\bigr)\\&\quad + 1\cdot (-1)^{1+2}\cdot\bigl((-2)\cdot 2-(-2)\cdot (-2)\bigr)\\&\quad + (-1)\cdot (-1)^{1+3}\cdot\bigl((-2)\cdot (-1)-(-2)\cdot (-2)\bigr)\\&=4\end{aligned}
\]
El \textbf{menor} es el número anterior. El \textbf{cofactor} incorpora además el signo $(-1)^{i+j}$:
\[
\boxed{C_{42}=(-1)^{4+2}M_{42}=1\cdot 4=4}.
\]

\par\bigskip\noindent\textbf{Solución Ejercicio 185}\par\nobreak
Consideremos $A=\begin{pmatrix}-1&-2&-2&1\\1&-1&1&1\\1&-2&-2&1\\0&-2&2&-1\end{pmatrix}$.

El elemento indicado ocupa la fila 2 y la columna 3. Primero \textbf{eliminamos esa fila y esa columna}.

\[
M_{23}=\begin{vmatrix}-1&-2&1\\1&-2&1\\0&-2&-1\end{vmatrix}.
\]
Calculamos este determinante de orden tres por cofactores:
\[
\begin{aligned}M_{23}&=(-2)\cdot (-1)^{3+2}\cdot\bigl((-1)\cdot 1-1\cdot 1\bigr)\\&\quad + (-1)\cdot (-1)^{3+3}\cdot\bigl((-1)\cdot (-2)-(-2)\cdot 1\bigr)\\&=-8\end{aligned}
\]
El \textbf{menor} es el número anterior. El \textbf{cofactor} incorpora además el signo $(-1)^{i+j}$:
\[
\boxed{C_{23}=(-1)^{2+3}M_{23}=(-1)\cdot (-8)=8}.
\]

\par\bigskip\noindent\textbf{Solución Ejercicio 187}\par\nobreak
Consideremos $A=\begin{pmatrix}1&0\\0&1\end{pmatrix}$.

Para obtener la adjunta seguimos el orden: \textbf{menores $\to$ cofactores $\to$ transposición $\to$ inversa}.

Comenzamos calculando el determinante de la matriz.
Usamos $\det\begin{psmallmatrix}a&b\\c&d\end{psmallmatrix}=ad-bc$:
\[
\begin{vmatrix}1&0\\0&1\end{vmatrix}=1\cdot 1-0\cdot 0=1-(0)=1.
\]
Calculamos ahora los menores y sus correspondientes cofactores. Calculamos $M_{ij}=\det(A_{ij})$ eliminando la fila $i$ y la columna $j$, y después $C_{ij}=(-1)^{i+j}M_{ij}$.

\[
M_{11}=1,\quad C_{11}=(-1)^{1+1}(1)=1.
\]
\[
M_{12}=0,\quad C_{12}=(-1)^{1+2}(0)=0.
\]
\[
M_{21}=0,\quad C_{21}=(-1)^{2+1}(0)=0.
\]
\[
M_{22}=1,\quad C_{22}=(-1)^{2+2}(1)=1.
\]
Reunimos los cofactores en una matriz y la transponemos para obtener la adjunta.
\[
\mathcal{C}(A)=\begin{pmatrix}1&0\\0&1\end{pmatrix},\qquad\adj(A)=\mathcal{C}(A)^T=\begin{pmatrix}1&0\\0&1\end{pmatrix}.
\]
Como el determinante es distinto de cero, obtenemos la inversa mediante la adjunta.
\[
A^{-1}=\frac1{\det(A)}\adj(A)=\frac1{1}\begin{pmatrix}1&0\\0&1\end{pmatrix}=\begin{pmatrix}1&0\\0&1\end{pmatrix}.
\]
Como verificación, debe cumplirse $A\adj(A)=\det(A)I$.
\[
\begin{pmatrix}1&0\\0&1\end{pmatrix}\begin{pmatrix}1&0\\0&1\end{pmatrix}=\begin{pmatrix}1&0\\0&1\end{pmatrix}=1I_{2},
\]
\[
AA^{-1}=\begin{pmatrix}1&0\\0&1\end{pmatrix}=I_{2}.
\]
Para comprobarlo también por Gauss--Jordan, reducimos $[A\mid I]$ hasta obtener $[I\mid A^{-1}]$:

Comenzamos con la matriz aumentada $[A\mid I]$ y aplicamos operaciones elementales; cada operación actúa sobre toda la fila.

\[
\left[\begin{array}{rr|rr}1&0&1&0\\0&1&0&1\end{array}\right]
\]
La matriz ya está en forma escalonada reducida; no se necesitan operaciones.

Como la parte izquierda es $I$, la parte derecha es la inversa:
\[
\boxed{A^{-1}=\begin{pmatrix}1&0\\0&1\end{pmatrix}}
\]
Como comprobación, multiplicamos y obtenemos
\[
AA^{-1}=\begin{pmatrix}1&0\\0&1\end{pmatrix}\begin{pmatrix}1&0\\0&1\end{pmatrix}=\begin{pmatrix}1&0\\0&1\end{pmatrix}=I_{2}.
\]

\par\bigskip\noindent\textbf{Solución Ejercicio 189}\par\nobreak
Consideremos $A=\begin{pmatrix}3&-1\\-2&1\end{pmatrix}$.

Para obtener la adjunta seguimos el orden: \textbf{menores $\to$ cofactores $\to$ transposición $\to$ inversa}.

Comenzamos calculando el determinante de la matriz.
Usamos $\det\begin{psmallmatrix}a&b\\c&d\end{psmallmatrix}=ad-bc$:
\[
\begin{vmatrix}3&-1\\-2&1\end{vmatrix}=3\cdot 1-(-1)\cdot (-2)=3-(2)=1.
\]
Calculamos ahora los menores y sus correspondientes cofactores. Calculamos $M_{ij}=\det(A_{ij})$ eliminando la fila $i$ y la columna $j$, y después $C_{ij}=(-1)^{i+j}M_{ij}$.

\[
M_{11}=1,\quad C_{11}=(-1)^{1+1}(1)=1.
\]
\[
M_{12}=-2,\quad C_{12}=(-1)^{1+2}(-2)=2.
\]
\[
M_{21}=-1,\quad C_{21}=(-1)^{2+1}(-1)=1.
\]
\[
M_{22}=3,\quad C_{22}=(-1)^{2+2}(3)=3.
\]
Reunimos los cofactores en una matriz y la transponemos para obtener la adjunta.
\[
\mathcal{C}(A)=\begin{pmatrix}1&2\\1&3\end{pmatrix},\qquad\adj(A)=\mathcal{C}(A)^T=\begin{pmatrix}1&1\\2&3\end{pmatrix}.
\]
Como el determinante es distinto de cero, obtenemos la inversa mediante la adjunta.
\[
A^{-1}=\frac1{\det(A)}\adj(A)=\frac1{1}\begin{pmatrix}1&1\\2&3\end{pmatrix}=\begin{pmatrix}1&1\\2&3\end{pmatrix}.
\]
Como verificación, debe cumplirse $A\adj(A)=\det(A)I$.
\[
\begin{pmatrix}3&-1\\-2&1\end{pmatrix}\begin{pmatrix}1&1\\2&3\end{pmatrix}=\begin{pmatrix}1&0\\0&1\end{pmatrix}=1I_{2},
\]
\[
AA^{-1}=\begin{pmatrix}1&0\\0&1\end{pmatrix}=I_{2}.
\]
Para comprobarlo también por Gauss--Jordan, reducimos $[A\mid I]$ hasta obtener $[I\mid A^{-1}]$:

Comenzamos con la matriz aumentada $[A\mid I]$ y aplicamos operaciones elementales; cada operación actúa sobre toda la fila.

\[
\left[\begin{array}{rr|rr}3&-1&1&0\\-2&1&0&1\end{array}\right]
\]
Operación $R_{1}\leftarrow \frac{1}{3}R_{1}$:
\[
\left[\begin{array}{rr|rr}1&- \frac{1}{3}&\frac{1}{3}&0\\-2&1&0&1\end{array}\right]
\]
Operación $R_{2}\leftarrow R_{2}+2R_{1}$:
\[
\left[\begin{array}{rr|rr}1&- \frac{1}{3}&\frac{1}{3}&0\\0&\frac{1}{3}&\frac{2}{3}&1\end{array}\right]
\]
Operación $R_{2}\leftarrow 3R_{2}$:
\[
\left[\begin{array}{rr|rr}1&- \frac{1}{3}&\frac{1}{3}&0\\0&1&2&3\end{array}\right]
\]
Operación $R_{1}\leftarrow R_{1}+\frac{1}{3}R_{2}$:
\[
\left[\begin{array}{rr|rr}1&0&1&1\\0&1&2&3\end{array}\right]
\]
Como la parte izquierda es $I$, la parte derecha es la inversa:
\[
\boxed{A^{-1}=\begin{pmatrix}1&1\\2&3\end{pmatrix}}
\]
Como comprobación, multiplicamos y obtenemos
\[
AA^{-1}=\begin{pmatrix}3&-1\\-2&1\end{pmatrix}\begin{pmatrix}1&1\\2&3\end{pmatrix}=\begin{pmatrix}1&0\\0&1\end{pmatrix}=I_{2}.
\]

\par\bigskip\noindent\textbf{Solución Ejercicio 191}\par\nobreak
Consideremos $A=\begin{pmatrix}3&2\\1&1\end{pmatrix}$.

Para obtener la adjunta seguimos el orden: \textbf{menores $\to$ cofactores $\to$ transposición $\to$ inversa}.

Comenzamos calculando el determinante de la matriz.
Usamos $\det\begin{psmallmatrix}a&b\\c&d\end{psmallmatrix}=ad-bc$:
\[
\begin{vmatrix}3&2\\1&1\end{vmatrix}=3\cdot 1-2\cdot 1=3-(2)=1.
\]
Calculamos ahora los menores y sus correspondientes cofactores. Calculamos $M_{ij}=\det(A_{ij})$ eliminando la fila $i$ y la columna $j$, y después $C_{ij}=(-1)^{i+j}M_{ij}$.

\[
M_{11}=1,\quad C_{11}=(-1)^{1+1}(1)=1.
\]
\[
M_{12}=1,\quad C_{12}=(-1)^{1+2}(1)=-1.
\]
\[
M_{21}=2,\quad C_{21}=(-1)^{2+1}(2)=-2.
\]
\[
M_{22}=3,\quad C_{22}=(-1)^{2+2}(3)=3.
\]
Reunimos los cofactores en una matriz y la transponemos para obtener la adjunta.
\[
\mathcal{C}(A)=\begin{pmatrix}1&-1\\-2&3\end{pmatrix},\qquad\adj(A)=\mathcal{C}(A)^T=\begin{pmatrix}1&-2\\-1&3\end{pmatrix}.
\]
Como el determinante es distinto de cero, obtenemos la inversa mediante la adjunta.
\[
A^{-1}=\frac1{\det(A)}\adj(A)=\frac1{1}\begin{pmatrix}1&-2\\-1&3\end{pmatrix}=\begin{pmatrix}1&-2\\-1&3\end{pmatrix}.
\]
Como verificación, debe cumplirse $A\adj(A)=\det(A)I$.
\[
\begin{pmatrix}3&2\\1&1\end{pmatrix}\begin{pmatrix}1&-2\\-1&3\end{pmatrix}=\begin{pmatrix}1&0\\0&1\end{pmatrix}=1I_{2},
\]
\[
AA^{-1}=\begin{pmatrix}1&0\\0&1\end{pmatrix}=I_{2}.
\]
Para comprobarlo también por Gauss--Jordan, reducimos $[A\mid I]$ hasta obtener $[I\mid A^{-1}]$:

Comenzamos con la matriz aumentada $[A\mid I]$ y aplicamos operaciones elementales; cada operación actúa sobre toda la fila.

\[
\left[\begin{array}{rr|rr}3&2&1&0\\1&1&0&1\end{array}\right]
\]
Operación $R_{1}\leftarrow \frac{1}{3}R_{1}$:
\[
\left[\begin{array}{rr|rr}1&\frac{2}{3}&\frac{1}{3}&0\\1&1&0&1\end{array}\right]
\]
Operación $R_{2}\leftarrow R_{2}-1R_{1}$:
\[
\left[\begin{array}{rr|rr}1&\frac{2}{3}&\frac{1}{3}&0\\0&\frac{1}{3}&- \frac{1}{3}&1\end{array}\right]
\]
Operación $R_{2}\leftarrow 3R_{2}$:
\[
\left[\begin{array}{rr|rr}1&\frac{2}{3}&\frac{1}{3}&0\\0&1&-1&3\end{array}\right]
\]
Operación $R_{1}\leftarrow R_{1}-\frac{2}{3}R_{2}$:
\[
\left[\begin{array}{rr|rr}1&0&1&-2\\0&1&-1&3\end{array}\right]
\]
Como la parte izquierda es $I$, la parte derecha es la inversa:
\[
\boxed{A^{-1}=\begin{pmatrix}1&-2\\-1&3\end{pmatrix}}
\]
Como comprobación, multiplicamos y obtenemos
\[
AA^{-1}=\begin{pmatrix}3&2\\1&1\end{pmatrix}\begin{pmatrix}1&-2\\-1&3\end{pmatrix}=\begin{pmatrix}1&0\\0&1\end{pmatrix}=I_{2}.
\]

\par\bigskip\noindent\textbf{Solución Ejercicio 193}\par\nobreak
Consideremos $A=\begin{pmatrix}1&0&0\\0&-1&3\\0&-1&2\end{pmatrix}$.

Para obtener la adjunta seguimos el orden: \textbf{menores $\to$ cofactores $\to$ transposición $\to$ inversa}.

Comenzamos calculando el determinante de la matriz.
Elegimos la fila 1, que contiene dos ceros, para reducir el trabajo. El patrón de signos es $C_{ij}=(-1)^{i+j}M_{ij}$, donde $M_{ij}$ es el determinante que queda al borrar la fila $i$ y la columna $j$.

Para la entrada $a_{11}=1$, borramos su fila y su columna:
\[
M_{11}=\begin{vmatrix}-1&3\\-1&2\end{vmatrix}.
\]
\[
M_{11}=(-1)\cdot 2-3\cdot (-1)=1.
\]
\[
C_{11}=(-1)^{1+1}(1)=1.
\]
Expandimos: \[
\boxed{\det(A)=1\cdot 1=1}
\]
Calculamos ahora los menores y sus correspondientes cofactores. Calculamos $M_{ij}=\det(A_{ij})$ eliminando la fila $i$ y la columna $j$, y después $C_{ij}=(-1)^{i+j}M_{ij}$.

Fila 1 de la matriz de cofactores:
\[
M_{11}=\begin{vmatrix}-1&3\\-1&2\end{vmatrix}=(-1)\cdot 2-3\cdot (-1)=1,\quad C_{11}=(-1)^{1+1}M_{11}=1.
\]
\[
M_{12}=\begin{vmatrix}0&3\\0&2\end{vmatrix}=0\cdot 2-3\cdot 0=0,\quad C_{12}=(-1)^{1+2}M_{12}=0.
\]
\[
M_{13}=\begin{vmatrix}0&-1\\0&-1\end{vmatrix}=0\cdot (-1)-(-1)\cdot 0=0,\quad C_{13}=(-1)^{1+3}M_{13}=0.
\]
Fila 2 de la matriz de cofactores:
\[
M_{21}=\begin{vmatrix}0&0\\-1&2\end{vmatrix}=0\cdot 2-0\cdot (-1)=0,\quad C_{21}=(-1)^{2+1}M_{21}=0.
\]
\[
M_{22}=\begin{vmatrix}1&0\\0&2\end{vmatrix}=1\cdot 2-0\cdot 0=2,\quad C_{22}=(-1)^{2+2}M_{22}=2.
\]
\[
M_{23}=\begin{vmatrix}1&0\\0&-1\end{vmatrix}=1\cdot (-1)-0\cdot 0=-1,\quad C_{23}=(-1)^{2+3}M_{23}=1.
\]
Fila 3 de la matriz de cofactores:
\[
M_{31}=\begin{vmatrix}0&0\\-1&3\end{vmatrix}=0\cdot 3-0\cdot (-1)=0,\quad C_{31}=(-1)^{3+1}M_{31}=0.
\]
\[
M_{32}=\begin{vmatrix}1&0\\0&3\end{vmatrix}=1\cdot 3-0\cdot 0=3,\quad C_{32}=(-1)^{3+2}M_{32}=-3.
\]
\[
M_{33}=\begin{vmatrix}1&0\\0&-1\end{vmatrix}=1\cdot (-1)-0\cdot 0=-1,\quad C_{33}=(-1)^{3+3}M_{33}=-1.
\]
Reunimos los cofactores en una matriz y la transponemos para obtener la adjunta.
\[
\mathcal{C}(A)=\begin{pmatrix}1&0&0\\0&2&1\\0&-3&-1\end{pmatrix},\qquad\adj(A)=\mathcal{C}(A)^T=\begin{pmatrix}1&0&0\\0&2&-3\\0&1&-1\end{pmatrix}.
\]
Como el determinante es distinto de cero, obtenemos la inversa mediante la adjunta.
\[
A^{-1}=\frac1{\det(A)}\adj(A)=\frac1{1}\begin{pmatrix}1&0&0\\0&2&-3\\0&1&-1\end{pmatrix}=\begin{pmatrix}1&0&0\\0&2&-3\\0&1&-1\end{pmatrix}.
\]
Como verificación, debe cumplirse $A\adj(A)=\det(A)I$.
\[
\begin{pmatrix}1&0&0\\0&-1&3\\0&-1&2\end{pmatrix}\begin{pmatrix}1&0&0\\0&2&-3\\0&1&-1\end{pmatrix}=\begin{pmatrix}1&0&0\\0&1&0\\0&0&1\end{pmatrix}=1I_{3},
\]
\[
AA^{-1}=\begin{pmatrix}1&0&0\\0&1&0\\0&0&1\end{pmatrix}=I_{3}.
\]
Para comprobarlo también por Gauss--Jordan, reducimos $[A\mid I]$ hasta obtener $[I\mid A^{-1}]$:

Comenzamos con la matriz aumentada $[A\mid I]$ y aplicamos operaciones elementales; cada operación actúa sobre toda la fila.

\[
\left[\begin{array}{rrr|rrr}1&0&0&1&0&0\\0&-1&3&0&1&0\\0&-1&2&0&0&1\end{array}\right]
\]
Operación $R_{2}\leftarrow -1R_{2}$:
\[
\left[\begin{array}{rrr|rrr}1&0&0&1&0&0\\0&1&-3&0&-1&0\\0&-1&2&0&0&1\end{array}\right]
\]
Operación $R_{3}\leftarrow R_{3}+1R_{2}$:
\[
\left[\begin{array}{rrr|rrr}1&0&0&1&0&0\\0&1&-3&0&-1&0\\0&0&-1&0&-1&1\end{array}\right]
\]
Operación $R_{3}\leftarrow -1R_{3}$:
\[
\left[\begin{array}{rrr|rrr}1&0&0&1&0&0\\0&1&-3&0&-1&0\\0&0&1&0&1&-1\end{array}\right]
\]
Operación $R_{2}\leftarrow R_{2}+3R_{3}$:
\[
\left[\begin{array}{rrr|rrr}1&0&0&1&0&0\\0&1&0&0&2&-3\\0&0&1&0&1&-1\end{array}\right]
\]
Como la parte izquierda es $I$, la parte derecha es la inversa:
\[
\boxed{A^{-1}=\begin{pmatrix}1&0&0\\0&2&-3\\0&1&-1\end{pmatrix}}
\]
Como comprobación, multiplicamos y obtenemos
\[
AA^{-1}=\begin{pmatrix}1&0&0\\0&-1&3\\0&-1&2\end{pmatrix}\begin{pmatrix}1&0&0\\0&2&-3\\0&1&-1\end{pmatrix}=\begin{pmatrix}1&0&0\\0&1&0\\0&0&1\end{pmatrix}=I_{3}.
\]

\par\bigskip\noindent\textbf{Solución Ejercicio 195}\par\nobreak
Consideremos $A=\begin{pmatrix}1&0&0\\-1&1&-2\\2&0&1\end{pmatrix}$.

Para obtener la adjunta seguimos el orden: \textbf{menores $\to$ cofactores $\to$ transposición $\to$ inversa}.

Comenzamos calculando el determinante de la matriz.
Elegimos la fila 1, que contiene dos ceros, para reducir el trabajo. El patrón de signos es $C_{ij}=(-1)^{i+j}M_{ij}$, donde $M_{ij}$ es el determinante que queda al borrar la fila $i$ y la columna $j$.

Para la entrada $a_{11}=1$, borramos su fila y su columna:
\[
M_{11}=\begin{vmatrix}1&-2\\0&1\end{vmatrix}.
\]
\[
M_{11}=1\cdot 1-(-2)\cdot 0=1.
\]
\[
C_{11}=(-1)^{1+1}(1)=1.
\]
Expandimos: \[
\boxed{\det(A)=1\cdot 1=1}
\]
Calculamos ahora los menores y sus correspondientes cofactores. Calculamos $M_{ij}=\det(A_{ij})$ eliminando la fila $i$ y la columna $j$, y después $C_{ij}=(-1)^{i+j}M_{ij}$.

Fila 1 de la matriz de cofactores:
\[
M_{11}=\begin{vmatrix}1&-2\\0&1\end{vmatrix}=1\cdot 1-(-2)\cdot 0=1,\quad C_{11}=(-1)^{1+1}M_{11}=1.
\]
\[
M_{12}=\begin{vmatrix}-1&-2\\2&1\end{vmatrix}=(-1)\cdot 1-(-2)\cdot 2=3,\quad C_{12}=(-1)^{1+2}M_{12}=-3.
\]
\[
M_{13}=\begin{vmatrix}-1&1\\2&0\end{vmatrix}=(-1)\cdot 0-1\cdot 2=-2,\quad C_{13}=(-1)^{1+3}M_{13}=-2.
\]
Fila 2 de la matriz de cofactores:
\[
M_{21}=\begin{vmatrix}0&0\\0&1\end{vmatrix}=0\cdot 1-0\cdot 0=0,\quad C_{21}=(-1)^{2+1}M_{21}=0.
\]
\[
M_{22}=\begin{vmatrix}1&0\\2&1\end{vmatrix}=1\cdot 1-0\cdot 2=1,\quad C_{22}=(-1)^{2+2}M_{22}=1.
\]
\[
M_{23}=\begin{vmatrix}1&0\\2&0\end{vmatrix}=1\cdot 0-0\cdot 2=0,\quad C_{23}=(-1)^{2+3}M_{23}=0.
\]
Fila 3 de la matriz de cofactores:
\[
M_{31}=\begin{vmatrix}0&0\\1&-2\end{vmatrix}=0\cdot (-2)-0\cdot 1=0,\quad C_{31}=(-1)^{3+1}M_{31}=0.
\]
\[
M_{32}=\begin{vmatrix}1&0\\-1&-2\end{vmatrix}=1\cdot (-2)-0\cdot (-1)=-2,\quad C_{32}=(-1)^{3+2}M_{32}=2.
\]
\[
M_{33}=\begin{vmatrix}1&0\\-1&1\end{vmatrix}=1\cdot 1-0\cdot (-1)=1,\quad C_{33}=(-1)^{3+3}M_{33}=1.
\]
Reunimos los cofactores en una matriz y la transponemos para obtener la adjunta.
\[
\mathcal{C}(A)=\begin{pmatrix}1&-3&-2\\0&1&0\\0&2&1\end{pmatrix},\qquad\adj(A)=\mathcal{C}(A)^T=\begin{pmatrix}1&0&0\\-3&1&2\\-2&0&1\end{pmatrix}.
\]
Como el determinante es distinto de cero, obtenemos la inversa mediante la adjunta.
\[
A^{-1}=\frac1{\det(A)}\adj(A)=\frac1{1}\begin{pmatrix}1&0&0\\-3&1&2\\-2&0&1\end{pmatrix}=\begin{pmatrix}1&0&0\\-3&1&2\\-2&0&1\end{pmatrix}.
\]
Como verificación, debe cumplirse $A\adj(A)=\det(A)I$.
\[
\begin{pmatrix}1&0&0\\-1&1&-2\\2&0&1\end{pmatrix}\begin{pmatrix}1&0&0\\-3&1&2\\-2&0&1\end{pmatrix}=\begin{pmatrix}1&0&0\\0&1&0\\0&0&1\end{pmatrix}=1I_{3},
\]
\[
AA^{-1}=\begin{pmatrix}1&0&0\\0&1&0\\0&0&1\end{pmatrix}=I_{3}.
\]
Para comprobarlo también por Gauss--Jordan, reducimos $[A\mid I]$ hasta obtener $[I\mid A^{-1}]$:

Comenzamos con la matriz aumentada $[A\mid I]$ y aplicamos operaciones elementales; cada operación actúa sobre toda la fila.

\[
\left[\begin{array}{rrr|rrr}1&0&0&1&0&0\\-1&1&-2&0&1&0\\2&0&1&0&0&1\end{array}\right]
\]
Operación $R_{2}\leftarrow R_{2}+1R_{1}$:
\[
\left[\begin{array}{rrr|rrr}1&0&0&1&0&0\\0&1&-2&1&1&0\\2&0&1&0&0&1\end{array}\right]
\]
Operación $R_{3}\leftarrow R_{3}-2R_{1}$:
\[
\left[\begin{array}{rrr|rrr}1&0&0&1&0&0\\0&1&-2&1&1&0\\0&0&1&-2&0&1\end{array}\right]
\]
Operación $R_{2}\leftarrow R_{2}+2R_{3}$:
\[
\left[\begin{array}{rrr|rrr}1&0&0&1&0&0\\0&1&0&-3&1&2\\0&0&1&-2&0&1\end{array}\right]
\]
Como la parte izquierda es $I$, la parte derecha es la inversa:
\[
\boxed{A^{-1}=\begin{pmatrix}1&0&0\\-3&1&2\\-2&0&1\end{pmatrix}}
\]
Como comprobación, multiplicamos y obtenemos
\[
AA^{-1}=\begin{pmatrix}1&0&0\\-1&1&-2\\2&0&1\end{pmatrix}\begin{pmatrix}1&0&0\\-3&1&2\\-2&0&1\end{pmatrix}=\begin{pmatrix}1&0&0\\0&1&0\\0&0&1\end{pmatrix}=I_{3}.
\]

\par\bigskip\noindent\textbf{Solución Ejercicio 197}\par\nobreak
Aplicamos operaciones elementales de fila, que no modifican el rango ni el conjunto de soluciones del sistema homogéneo.

\[
\begin{pmatrix}2&0&2\\-1&0&-1\end{pmatrix}
\]
Operación $R_{1}\leftarrow \frac{1}{2}R_{1}$:
\[
\begin{pmatrix}1&0&1\\-1&0&-1\end{pmatrix}
\]
Operación $R_{2}\leftarrow R_{2}+1R_{1}$:
\[
\begin{pmatrix}1&0&1\\0&0&0\end{pmatrix}
\]
Hay pivotes en las columnas 1; por tanto, $\boxed{\rang(A)=1}$.

En la matriz reducida contamos 1 fila(s) no nula(s), una por cada pivote.


\par\bigskip\noindent\textbf{Solución Ejercicio 199}\par\nobreak
Aplicamos operaciones elementales de fila, que no modifican el rango ni el conjunto de soluciones del sistema homogéneo.

\[
\begin{pmatrix}1&2\\0&0\\0&0\end{pmatrix}
\]
La matriz ya está en forma escalonada reducida; no se necesitan operaciones.

Hay pivotes en las columnas 1; por tanto, $\boxed{\rang(A)=1}$.

En la matriz reducida contamos 1 fila(s) no nula(s), una por cada pivote.


\par\bigskip\noindent\textbf{Solución Ejercicio 201}\par\nobreak
Aplicamos operaciones elementales de fila, que no modifican el rango ni el conjunto de soluciones del sistema homogéneo.

\[
\begin{pmatrix}1&0&0\\0&1&0\\0&1&1\\0&-1&-1\end{pmatrix}
\]
Operación $R_{3}\leftarrow R_{3}-1R_{2}$:
\[
\begin{pmatrix}1&0&0\\0&1&0\\0&0&1\\0&-1&-1\end{pmatrix}
\]
Operación $R_{4}\leftarrow R_{4}+1R_{2}$:
\[
\begin{pmatrix}1&0&0\\0&1&0\\0&0&1\\0&0&-1\end{pmatrix}
\]
Operación $R_{4}\leftarrow R_{4}+1R_{3}$:
\[
\begin{pmatrix}1&0&0\\0&1&0\\0&0&1\\0&0&0\end{pmatrix}
\]
Hay pivotes en las columnas 1, 2, 3; por tanto, $\boxed{\rang(A)=3}$.

En la matriz reducida contamos 3 fila(s) no nula(s), una por cada pivote.


\par\bigskip\noindent\textbf{Solución Ejercicio 203}\par\nobreak
Aplicamos operaciones elementales de fila, que no modifican el rango ni el conjunto de soluciones del sistema homogéneo.

\[
\begin{pmatrix}1&0&0&1\\0&1&0&-1\\0&1&1&0\\0&1&1&0\\0&0&0&0\end{pmatrix}
\]
Operación $R_{3}\leftarrow R_{3}-1R_{2}$:
\[
\begin{pmatrix}1&0&0&1\\0&1&0&-1\\0&0&1&1\\0&1&1&0\\0&0&0&0\end{pmatrix}
\]
Operación $R_{4}\leftarrow R_{4}-1R_{2}$:
\[
\begin{pmatrix}1&0&0&1\\0&1&0&-1\\0&0&1&1\\0&0&1&1\\0&0&0&0\end{pmatrix}
\]
Operación $R_{4}\leftarrow R_{4}-1R_{3}$:
\[
\begin{pmatrix}1&0&0&1\\0&1&0&-1\\0&0&1&1\\0&0&0&0\\0&0&0&0\end{pmatrix}
\]
Hay pivotes en las columnas 1, 2, 3; por tanto, $\boxed{\rang(A)=3}$.

En la matriz reducida contamos 3 fila(s) no nula(s), una por cada pivote.


\par\bigskip\noindent\textbf{Solución Ejercicio 205}\par\nobreak
Aplicamos operaciones elementales de fila, que no modifican el rango ni el conjunto de soluciones del sistema homogéneo.

\[
\begin{pmatrix}1&0\\0&1\\0&0\\0&0\\0&0\end{pmatrix}
\]
La matriz ya está en forma escalonada reducida; no se necesitan operaciones.

Hay pivotes en las columnas 1, 2; por tanto, $\boxed{\rang(A)=2}$.

En la matriz reducida contamos 2 fila(s) no nula(s), una por cada pivote.


\par\bigskip\noindent\textbf{Solución Ejercicio 207}\par\nobreak
Aplicamos operaciones elementales de fila, que no modifican el rango ni el conjunto de soluciones del sistema homogéneo.

\[
\begin{pmatrix}1&1&0&3\\0&1&0&1\\0&1&1&2\\0&0&0&0\end{pmatrix}
\]
Operación $R_{1}\leftarrow R_{1}-1R_{2}$:
\[
\begin{pmatrix}1&0&0&2\\0&1&0&1\\0&1&1&2\\0&0&0&0\end{pmatrix}
\]
Operación $R_{3}\leftarrow R_{3}-1R_{2}$:
\[
\begin{pmatrix}1&0&0&2\\0&1&0&1\\0&0&1&1\\0&0&0&0\end{pmatrix}
\]
Hay pivotes en las columnas 1, 2, 3; por tanto, $\boxed{\rang(A)=3}$.

En la matriz reducida contamos 3 fila(s) no nula(s), una por cada pivote.


\par\bigskip\noindent\textbf{Solución Ejercicio 209}\par\nobreak
Aplicamos operaciones elementales de fila, que no modifican el rango ni el conjunto de soluciones del sistema homogéneo.

\[
\begin{pmatrix}1&2&0\\0&0&0\\0&0&0\\0&0&0\\0&0&0\end{pmatrix}
\]
La matriz ya está en forma escalonada reducida; no se necesitan operaciones.

Hay pivotes en las columnas 1; por tanto, $\boxed{\rang(A)=1}$.

En la matriz reducida contamos 1 fila(s) no nula(s), una por cada pivote.


\par\bigskip\noindent\textbf{Solución Ejercicio 211}\par\nobreak
Aplicamos operaciones elementales de fila, que no modifican el rango ni el conjunto de soluciones del sistema homogéneo.

\[
\begin{pmatrix}1&-1&0&0&-1\\0&1&0&0&2\\0&0&1&0&-1\\-1&0&0&1&-1\end{pmatrix}
\]
Operación $R_{4}\leftarrow R_{4}+1R_{1}$:
\[
\begin{pmatrix}1&-1&0&0&-1\\0&1&0&0&2\\0&0&1&0&-1\\0&-1&0&1&-2\end{pmatrix}
\]
Operación $R_{1}\leftarrow R_{1}+1R_{2}$:
\[
\begin{pmatrix}1&0&0&0&1\\0&1&0&0&2\\0&0&1&0&-1\\0&-1&0&1&-2\end{pmatrix}
\]
Operación $R_{4}\leftarrow R_{4}+1R_{2}$:
\[
\begin{pmatrix}1&0&0&0&1\\0&1&0&0&2\\0&0&1&0&-1\\0&0&0&1&0\end{pmatrix}
\]
Hay pivotes en las columnas 1, 2, 3, 4; por tanto, $\boxed{\rang(A)=4}$.

En la matriz reducida contamos 4 fila(s) no nula(s), una por cada pivote.


\par\bigskip\noindent\textbf{Solución Ejercicio 213}\par\nobreak
Aplicamos operaciones elementales de fila, que no modifican el rango ni el conjunto de soluciones del sistema homogéneo.

\[
\begin{pmatrix}1&2&0\\-2&-4&0\end{pmatrix}
\]
Operación $R_{2}\leftarrow R_{2}+2R_{1}$:
\[
\begin{pmatrix}1&2&0\\0&0&0\end{pmatrix}
\]
Hay pivotes en las columnas 1; por tanto, $\boxed{\rang(A)=1}$.

Resolvemos $A\vect x=\vect0$; la matriz reducida representa el mismo sistema homogéneo.

Las variables libres son $x_{2}$, $x_{3}$. Sean $x_{2}=t_{1}$, $x_{3}=t_{2}$, con $t_i\in\mathbb R$.

\[
\begin{aligned}x_{1}=- 2 t_{1}\end{aligned}
\]
\[
\vect x=t_{1}\begin{pmatrix}-2\\1\\0\end{pmatrix} + t_{2}\begin{pmatrix}0\\0\\1\end{pmatrix}
\]
\[
\boxed{\Ker(A)=\operatorname{span}\left\{\begin{pmatrix}-2\\1\\0\end{pmatrix},\begin{pmatrix}0\\0\\1\end{pmatrix}\right\}}
\]
Para verificar el resultado, cada vector generador produce el vector cero al multiplicarlo por $A$:
\[
A\begin{pmatrix}-2\\1\\0\end{pmatrix}=\begin{pmatrix}0\\0\end{pmatrix}.
\]
\[
A\begin{pmatrix}0\\0\\1\end{pmatrix}=\begin{pmatrix}0\\0\end{pmatrix}.
\]
Además, $\rang(A)+\dim\Ker(A)=1+2=3$, igual al número de columnas.


\par\bigskip\noindent\textbf{Solución Ejercicio 215}\par\nobreak
Aplicamos operaciones elementales de fila, que no modifican el rango ni el conjunto de soluciones del sistema homogéneo.

\[
\begin{pmatrix}1&0\\0&1\\-2&0\end{pmatrix}
\]
Operación $R_{3}\leftarrow R_{3}+2R_{1}$:
\[
\begin{pmatrix}1&0\\0&1\\0&0\end{pmatrix}
\]
Hay pivotes en las columnas 1, 2; por tanto, $\boxed{\rang(A)=2}$.

Resolvemos $A\vect x=\vect0$; la matriz reducida representa el mismo sistema homogéneo.

\[
\begin{aligned}x_{1}=0\\x_{2}=0\end{aligned}
\]
No hay variables libres; la única solución es el vector nulo.
\[
\boxed{\Ker(A)=\{\vect0\}}
\]
Además, $\rang(A)+\dim\Ker(A)=2+0=2$, igual al número de columnas.


\par\bigskip\noindent\textbf{Solución Ejercicio 217}\par\nobreak
Aplicamos operaciones elementales de fila, que no modifican el rango ni el conjunto de soluciones del sistema homogéneo.

\[
\begin{pmatrix}1&0&0\\0&1&0\\-1&0&1\\0&-1&0\end{pmatrix}
\]
Operación $R_{3}\leftarrow R_{3}+1R_{1}$:
\[
\begin{pmatrix}1&0&0\\0&1&0\\0&0&1\\0&-1&0\end{pmatrix}
\]
Operación $R_{4}\leftarrow R_{4}+1R_{2}$:
\[
\begin{pmatrix}1&0&0\\0&1&0\\0&0&1\\0&0&0\end{pmatrix}
\]
Hay pivotes en las columnas 1, 2, 3; por tanto, $\boxed{\rang(A)=3}$.

Resolvemos $A\vect x=\vect0$; la matriz reducida representa el mismo sistema homogéneo.

\[
\begin{aligned}x_{1}=0\\x_{2}=0\\x_{3}=0\end{aligned}
\]
No hay variables libres; la única solución es el vector nulo.
\[
\boxed{\Ker(A)=\{\vect0\}}
\]
Además, $\rang(A)+\dim\Ker(A)=3+0=3$, igual al número de columnas.


\par\bigskip\noindent\textbf{Solución Ejercicio 219}\par\nobreak
Aplicamos operaciones elementales de fila, que no modifican el rango ni el conjunto de soluciones del sistema homogéneo.

\[
\begin{pmatrix}1&1&0&1\\0&1&0&1\\0&0&1&-1\\0&0&0&0\\1&1&0&1\end{pmatrix}
\]
Operación $R_{5}\leftarrow R_{5}-1R_{1}$:
\[
\begin{pmatrix}1&1&0&1\\0&1&0&1\\0&0&1&-1\\0&0&0&0\\0&0&0&0\end{pmatrix}
\]
Operación $R_{1}\leftarrow R_{1}-1R_{2}$:
\[
\begin{pmatrix}1&0&0&0\\0&1&0&1\\0&0&1&-1\\0&0&0&0\\0&0&0&0\end{pmatrix}
\]
Hay pivotes en las columnas 1, 2, 3; por tanto, $\boxed{\rang(A)=3}$.

Resolvemos $A\vect x=\vect0$; la matriz reducida representa el mismo sistema homogéneo.

Las variables libres son $x_{4}$. Sean $x_{4}=t_{1}$, con $t_i\in\mathbb R$.

\[
\begin{aligned}x_{1}=0\\x_{2}=- t_{1}\\x_{3}=t_{1}\end{aligned}
\]
\[
\vect x=t_{1}\begin{pmatrix}0\\-1\\1\\1\end{pmatrix}
\]
\[
\boxed{\Ker(A)=\operatorname{span}\left\{\begin{pmatrix}0\\-1\\1\\1\end{pmatrix}\right\}}
\]
Para verificar el resultado, cada vector generador produce el vector cero al multiplicarlo por $A$:
\[
A\begin{pmatrix}0\\-1\\1\\1\end{pmatrix}=\begin{pmatrix}0\\0\\0\\0\\0\end{pmatrix}.
\]
Además, $\rang(A)+\dim\Ker(A)=3+1=4$, igual al número de columnas.


\par\bigskip\noindent\textbf{Solución Ejercicio 221}\par\nobreak
Aplicamos operaciones elementales de fila, que no modifican el rango ni el conjunto de soluciones del sistema homogéneo.

\[
\begin{pmatrix}1&0\\1&1\\1&1\\0&0\\0&0\end{pmatrix}
\]
Operación $R_{2}\leftarrow R_{2}-1R_{1}$:
\[
\begin{pmatrix}1&0\\0&1\\1&1\\0&0\\0&0\end{pmatrix}
\]
Operación $R_{3}\leftarrow R_{3}-1R_{1}$:
\[
\begin{pmatrix}1&0\\0&1\\0&1\\0&0\\0&0\end{pmatrix}
\]
Operación $R_{3}\leftarrow R_{3}-1R_{2}$:
\[
\begin{pmatrix}1&0\\0&1\\0&0\\0&0\\0&0\end{pmatrix}
\]
Hay pivotes en las columnas 1, 2; por tanto, $\boxed{\rang(A)=2}$.

Resolvemos $A\vect x=\vect0$; la matriz reducida representa el mismo sistema homogéneo.

\[
\begin{aligned}x_{1}=0\\x_{2}=0\end{aligned}
\]
No hay variables libres; la única solución es el vector nulo.
\[
\boxed{\Ker(A)=\{\vect0\}}
\]
Además, $\rang(A)+\dim\Ker(A)=2+0=2$, igual al número de columnas.


\par\bigskip\noindent\textbf{Solución Ejercicio 223}\par\nobreak
Aplicamos operaciones elementales de fila, que no modifican el rango ni el conjunto de soluciones del sistema homogéneo.

\[
\begin{pmatrix}1&1&0&1\\0&1&0&-1\\0&0&1&2\\0&-1&0&1\end{pmatrix}
\]
Operación $R_{1}\leftarrow R_{1}-1R_{2}$:
\[
\begin{pmatrix}1&0&0&2\\0&1&0&-1\\0&0&1&2\\0&-1&0&1\end{pmatrix}
\]
Operación $R_{4}\leftarrow R_{4}+1R_{2}$:
\[
\begin{pmatrix}1&0&0&2\\0&1&0&-1\\0&0&1&2\\0&0&0&0\end{pmatrix}
\]
Hay pivotes en las columnas 1, 2, 3; por tanto, $\boxed{\rang(A)=3}$.

Resolvemos $A\vect x=\vect0$; la matriz reducida representa el mismo sistema homogéneo.

Las variables libres son $x_{4}$. Sean $x_{4}=t_{1}$, con $t_i\in\mathbb R$.

\[
\begin{aligned}x_{1}=- 2 t_{1}\\x_{2}=t_{1}\\x_{3}=- 2 t_{1}\end{aligned}
\]
\[
\vect x=t_{1}\begin{pmatrix}-2\\1\\-2\\1\end{pmatrix}
\]
\[
\boxed{\Ker(A)=\operatorname{span}\left\{\begin{pmatrix}-2\\1\\-2\\1\end{pmatrix}\right\}}
\]
Para verificar el resultado, cada vector generador produce el vector cero al multiplicarlo por $A$:
\[
A\begin{pmatrix}-2\\1\\-2\\1\end{pmatrix}=\begin{pmatrix}0\\0\\0\\0\end{pmatrix}.
\]
Además, $\rang(A)+\dim\Ker(A)=3+1=4$, igual al número de columnas.


\par\bigskip\noindent\textbf{Solución Ejercicio 225}\par\nobreak
Aplicamos operaciones elementales de fila, que no modifican el rango ni el conjunto de soluciones del sistema homogéneo.

\[
\begin{pmatrix}1&-2&-1\\0&0&0\\1&-2&-1\\0&0&0\\0&0&0\end{pmatrix}
\]
Operación $R_{3}\leftarrow R_{3}-1R_{1}$:
\[
\begin{pmatrix}1&-2&-1\\0&0&0\\0&0&0\\0&0&0\\0&0&0\end{pmatrix}
\]
Hay pivotes en las columnas 1; por tanto, $\boxed{\rang(A)=1}$.

Resolvemos $A\vect x=\vect0$; la matriz reducida representa el mismo sistema homogéneo.

Las variables libres son $x_{2}$, $x_{3}$. Sean $x_{2}=t_{1}$, $x_{3}=t_{2}$, con $t_i\in\mathbb R$.

\[
\begin{aligned}x_{1}=2 t_{1} + t_{2}\end{aligned}
\]
\[
\vect x=t_{1}\begin{pmatrix}2\\1\\0\end{pmatrix} + t_{2}\begin{pmatrix}1\\0\\1\end{pmatrix}
\]
\[
\boxed{\Ker(A)=\operatorname{span}\left\{\begin{pmatrix}2\\1\\0\end{pmatrix},\begin{pmatrix}1\\0\\1\end{pmatrix}\right\}}
\]
Para verificar el resultado, cada vector generador produce el vector cero al multiplicarlo por $A$:
\[
A\begin{pmatrix}2\\1\\0\end{pmatrix}=\begin{pmatrix}0\\0\\0\\0\\0\end{pmatrix}.
\]
\[
A\begin{pmatrix}1\\0\\1\end{pmatrix}=\begin{pmatrix}0\\0\\0\\0\\0\end{pmatrix}.
\]
Además, $\rang(A)+\dim\Ker(A)=1+2=3$, igual al número de columnas.


\par\bigskip\noindent\textbf{Solución Ejercicio 227}\par\nobreak
Aplicamos operaciones elementales de fila, que no modifican el rango ni el conjunto de soluciones del sistema homogéneo.

\[
\begin{pmatrix}1&0&0&0&2\\0&1&0&-2&3\\0&0&1&0&2\\0&0&0&1&-2\end{pmatrix}
\]
Operación $R_{2}\leftarrow R_{2}+2R_{4}$:
\[
\begin{pmatrix}1&0&0&0&2\\0&1&0&0&-1\\0&0&1&0&2\\0&0&0&1&-2\end{pmatrix}
\]
Hay pivotes en las columnas 1, 2, 3, 4; por tanto, $\boxed{\rang(A)=4}$.

Resolvemos $A\vect x=\vect0$; la matriz reducida representa el mismo sistema homogéneo.

Las variables libres son $x_{5}$. Sean $x_{5}=t_{1}$, con $t_i\in\mathbb R$.

\[
\begin{aligned}x_{1}=- 2 t_{1}\\x_{2}=t_{1}\\x_{3}=- 2 t_{1}\\x_{4}=2 t_{1}\end{aligned}
\]
\[
\vect x=t_{1}\begin{pmatrix}-2\\1\\-2\\2\\1\end{pmatrix}
\]
\[
\boxed{\Ker(A)=\operatorname{span}\left\{\begin{pmatrix}-2\\1\\-2\\2\\1\end{pmatrix}\right\}}
\]
Para verificar el resultado, cada vector generador produce el vector cero al multiplicarlo por $A$:
\[
A\begin{pmatrix}-2\\1\\-2\\2\\1\end{pmatrix}=\begin{pmatrix}0\\0\\0\\0\end{pmatrix}.
\]
Además, $\rang(A)+\dim\Ker(A)=4+1=5$, igual al número de columnas.


\medskip
